\documentclass[11pt,a4paper]{article}
% =====================================================================
%  Preámbulo común para los apuntes de Física Electrónica
%  Uso (desde unidades/XX-nombre/archivo.tex):
%     \documentclass[11pt,a4paper]{article}
%     % =====================================================================
%  Preámbulo común para los apuntes de Física Electrónica
%  Uso (desde unidades/XX-nombre/archivo.tex):
%     \documentclass[11pt,a4paper]{article}
%     % =====================================================================
%  Preámbulo común para los apuntes de Física Electrónica
%  Uso (desde unidades/XX-nombre/archivo.tex):
%     \documentclass[11pt,a4paper]{article}
%     \input{../../comun/preambulo.tex}
% =====================================================================

% ---------- Idioma y tipografía ----------
\usepackage[T1]{fontenc}
\usepackage[utf8]{inputenc}
\usepackage[spanish,es-nodecimaldot,es-noshorthands,es-tabla]{babel}
\usepackage{lmodern}
\usepackage{microtype}
\usepackage[a4paper,margin=2.3cm,headheight=14pt]{geometry}
\usepackage{parskip}
\setlength{\parskip}{0.55em}

% ---------- Matemática y unidades ----------
\usepackage{amsmath,amssymb,mathtools}
\usepackage{bm}
\usepackage[per-mode=symbol,output-decimal-marker={.},separate-uncertainty]{siunitx}
\sisetup{range-phrase={\text{--}},range-units=single}
\DeclareSIUnit{\eV}{eV}

% ---------- Gráficos ----------
\usepackage{graphicx}
\usepackage{xcolor}
\usepackage{tikz}
\usetikzlibrary{arrows.meta,decorations.pathmorphing,decorations.markings,
  patterns,calc,positioning,shapes.geometric,fit,backgrounds,shadings}
\usepackage{pgfplots}
\pgfplotsset{compat=1.18}
\usepgfplotslibrary{fillbetween}
\usepackage[siunitx,RPvoltages]{circuitikz}

% ---------- Tablas, listas, figuras ----------
\usepackage{booktabs,array,tabularx,multirow}
\usepackage{enumitem}
\setlist{itemsep=0.2em,topsep=0.3em}
\usepackage{caption,subcaption}
\captionsetup{font=small,labelfont=bf,format=hang}
\usepackage{float}
\usepackage{wrapfig}
\usepackage{needspace}

% ---------- Colores ----------
\definecolor{azul}{HTML}{1F4E79}
\definecolor{azulclaro}{HTML}{DCE9F5}
\definecolor{verde}{HTML}{2E7D32}
\definecolor{verdeclaro}{HTML}{E3F2E1}
\definecolor{naranja}{HTML}{C55A11}
\definecolor{naranjaclaro}{HTML}{FBE9DB}
\definecolor{rojo}{HTML}{B71C1C}
\definecolor{rojoclaro}{HTML}{FDECEA}
\definecolor{violeta}{HTML}{5E35B1}
\definecolor{violetaclaro}{HTML}{EDE7F6}
\definecolor{gris}{HTML}{555555}
% colores para dibujos de bandas / portadores
\colorlet{colEc}{azul}
\colorlet{colEv}{rojo}
\colorlet{colElec}{azul}
\colorlet{colHueco}{rojo}
\colorlet{colFoton}{naranja}

% ---------- Cajas didácticas ----------
\usepackage[most]{tcolorbox}
\tcbset{
  base/.style={enhanced,breakable,arc=2pt,boxrule=0.6pt,left=6pt,right=6pt,
    top=4pt,bottom=4pt,fonttitle=\bfseries\small,before skip=8pt,after skip=8pt}
}
% Definición
\newtcolorbox{definicion}[1][]{base,colback=azulclaro,colframe=azul,
  title={Definición\ifx\relax#1\relax\else: #1\fi}}
% Resultado / fórmula clave
\newtcolorbox{clave}[1][]{base,colback=verdeclaro,colframe=verde,
  title={\ifx\relax#1\relax Resultado clave\else #1\fi}}
% Ejemplo resuelto
\newtcolorbox{ejemplo}[1][]{base,colback=white,colframe=violeta,
  coltitle=white,colbacktitle=violeta,
  title={Ejemplo\ifx\relax#1\relax\else: #1\fi}}
% Corrección / aclaración respecto de las diapositivas
\newtcolorbox{ojo}[1][]{base,colback=rojoclaro,colframe=rojo,
  title={\ensuremath{\triangle}\ Ojo\ifx\relax#1\relax\else: #1\fi}}
% Intuición física / comentario
\newtcolorbox{intuicion}[1][]{base,colback=naranjaclaro,colframe=naranja,
  title={Idea física\ifx\relax#1\relax\else: #1\fi}}

% ---------- Encabezados ----------
\usepackage{fancyhdr}
\pagestyle{fancy}
\fancyhf{}
\renewcommand{\headrulewidth}{0.4pt}
\fancyhead[L]{\small\color{gris}Física Electrónica}
\fancyhead[R]{\small\color{gris}\unidadtitulo}
\fancyfoot[C]{\small\thepage}
\newcommand{\unidadtitulo}{}
\newcommand{\unidad}[1]{\renewcommand{\unidadtitulo}{#1}}

\usepackage{titlesec}
\titleformat{\section}{\Large\bfseries\color{azul}}{\thesection}{0.6em}{}
\titleformat{\subsection}{\large\bfseries\color{azul}}{\thesubsection}{0.6em}{}
\titleformat{\subsubsection}{\normalsize\bfseries\color{azul!80!black}}{\thesubsubsection}{0.6em}{}

\usepackage[hidelinks,bookmarksnumbered]{hyperref}
\usepackage{bookmark}

% ---------- Macros físicas ----------
\newcommand{\EG}{E_G}
\newcommand{\Ec}{E_c}
\newcommand{\Ev}{E_v}
\newcommand{\EF}{E_F}
\newcommand{\Eph}{E_{\mathrm{ph}}}
\newcommand{\lG}{\lambda_G}
\newcommand{\VT}{V_T}
\newcommand{\Popt}{P_{\mathrm{opt}}}
\newcommand{\Resp}{\mathcal{R}}
\newcommand{\vsat}{v_{\mathrm{sat}}}
\newcommand{\Voc}{V_{oc}}
\newcommand{\Isc}{I_{sc}}
\newcommand{\diapo}[1]{{\small\textsf{[diapo~#1]}}}
\newcommand{\dd}{\,\mathrm{d}}

% Distintivo para resultados que ya están en la hoja de fórmulas oficial.
% Uso: en el título de una caja, \begin{clave}[Título \enhoja], o en texto con \enhojatexto.
\newcommand{\enhojatexto}{{\scriptsize\sffamily\colorbox{verdeclaro}{\textcolor{verde!70!black}{\ensuremath{\checkmark}\,en la hoja}}}}
\newcommand{\enhoja}{\hfill{\scriptsize\sffamily\colorbox{white}{\textcolor{verde!70!black}{\ensuremath{\checkmark}\,en la hoja}}}}

\input{../../comun/constantes.tex}

% ---------- Estilos TikZ comunes ----------
\tikzset{
  banda/.style={line width=1.1pt},
  bandac/.style={banda,colEc},
  bandav/.style={banda,colEv},
  fermi/.style={dashed,gris,line width=0.6pt},
  elec/.style={circle,fill=colElec,inner sep=0pt,minimum size=5pt},
  hueco/.style={circle,draw=colHueco,fill=white,line width=0.9pt,inner sep=0pt,minimum size=5pt},
  foton/.style={colFoton,line width=0.9pt,decorate,
    decoration={snake,amplitude=1.6pt,segment length=5pt,post length=3pt},-{Stealth[length=5pt]}},
  flecha/.style={-{Stealth[length=5pt]},line width=0.7pt},
}

% =====================================================================

% ---------- Idioma y tipografía ----------
\usepackage[T1]{fontenc}
\usepackage[utf8]{inputenc}
\usepackage[spanish,es-nodecimaldot,es-noshorthands,es-tabla]{babel}
\usepackage{lmodern}
\usepackage{microtype}
\usepackage[a4paper,margin=2.3cm,headheight=14pt]{geometry}
\usepackage{parskip}
\setlength{\parskip}{0.55em}

% ---------- Matemática y unidades ----------
\usepackage{amsmath,amssymb,mathtools}
\usepackage{bm}
\usepackage[per-mode=symbol,output-decimal-marker={.},separate-uncertainty]{siunitx}
\sisetup{range-phrase={\text{--}},range-units=single}
\DeclareSIUnit{\eV}{eV}

% ---------- Gráficos ----------
\usepackage{graphicx}
\usepackage{xcolor}
\usepackage{tikz}
\usetikzlibrary{arrows.meta,decorations.pathmorphing,decorations.markings,
  patterns,calc,positioning,shapes.geometric,fit,backgrounds,shadings}
\usepackage{pgfplots}
\pgfplotsset{compat=1.18}
\usepgfplotslibrary{fillbetween}
\usepackage[siunitx,RPvoltages]{circuitikz}

% ---------- Tablas, listas, figuras ----------
\usepackage{booktabs,array,tabularx,multirow}
\usepackage{enumitem}
\setlist{itemsep=0.2em,topsep=0.3em}
\usepackage{caption,subcaption}
\captionsetup{font=small,labelfont=bf,format=hang}
\usepackage{float}
\usepackage{wrapfig}
\usepackage{needspace}

% ---------- Colores ----------
\definecolor{azul}{HTML}{1F4E79}
\definecolor{azulclaro}{HTML}{DCE9F5}
\definecolor{verde}{HTML}{2E7D32}
\definecolor{verdeclaro}{HTML}{E3F2E1}
\definecolor{naranja}{HTML}{C55A11}
\definecolor{naranjaclaro}{HTML}{FBE9DB}
\definecolor{rojo}{HTML}{B71C1C}
\definecolor{rojoclaro}{HTML}{FDECEA}
\definecolor{violeta}{HTML}{5E35B1}
\definecolor{violetaclaro}{HTML}{EDE7F6}
\definecolor{gris}{HTML}{555555}
% colores para dibujos de bandas / portadores
\colorlet{colEc}{azul}
\colorlet{colEv}{rojo}
\colorlet{colElec}{azul}
\colorlet{colHueco}{rojo}
\colorlet{colFoton}{naranja}

% ---------- Cajas didácticas ----------
\usepackage[most]{tcolorbox}
\tcbset{
  base/.style={enhanced,breakable,arc=2pt,boxrule=0.6pt,left=6pt,right=6pt,
    top=4pt,bottom=4pt,fonttitle=\bfseries\small,before skip=8pt,after skip=8pt}
}
% Definición
\newtcolorbox{definicion}[1][]{base,colback=azulclaro,colframe=azul,
  title={Definición\ifx\relax#1\relax\else: #1\fi}}
% Resultado / fórmula clave
\newtcolorbox{clave}[1][]{base,colback=verdeclaro,colframe=verde,
  title={\ifx\relax#1\relax Resultado clave\else #1\fi}}
% Ejemplo resuelto
\newtcolorbox{ejemplo}[1][]{base,colback=white,colframe=violeta,
  coltitle=white,colbacktitle=violeta,
  title={Ejemplo\ifx\relax#1\relax\else: #1\fi}}
% Corrección / aclaración respecto de las diapositivas
\newtcolorbox{ojo}[1][]{base,colback=rojoclaro,colframe=rojo,
  title={\ensuremath{\triangle}\ Ojo\ifx\relax#1\relax\else: #1\fi}}
% Intuición física / comentario
\newtcolorbox{intuicion}[1][]{base,colback=naranjaclaro,colframe=naranja,
  title={Idea física\ifx\relax#1\relax\else: #1\fi}}

% ---------- Encabezados ----------
\usepackage{fancyhdr}
\pagestyle{fancy}
\fancyhf{}
\renewcommand{\headrulewidth}{0.4pt}
\fancyhead[L]{\small\color{gris}Física Electrónica}
\fancyhead[R]{\small\color{gris}\unidadtitulo}
\fancyfoot[C]{\small\thepage}
\newcommand{\unidadtitulo}{}
\newcommand{\unidad}[1]{\renewcommand{\unidadtitulo}{#1}}

\usepackage{titlesec}
\titleformat{\section}{\Large\bfseries\color{azul}}{\thesection}{0.6em}{}
\titleformat{\subsection}{\large\bfseries\color{azul}}{\thesubsection}{0.6em}{}
\titleformat{\subsubsection}{\normalsize\bfseries\color{azul!80!black}}{\thesubsubsection}{0.6em}{}

\usepackage[hidelinks,bookmarksnumbered]{hyperref}
\usepackage{bookmark}

% ---------- Macros físicas ----------
\newcommand{\EG}{E_G}
\newcommand{\Ec}{E_c}
\newcommand{\Ev}{E_v}
\newcommand{\EF}{E_F}
\newcommand{\Eph}{E_{\mathrm{ph}}}
\newcommand{\lG}{\lambda_G}
\newcommand{\VT}{V_T}
\newcommand{\Popt}{P_{\mathrm{opt}}}
\newcommand{\Resp}{\mathcal{R}}
\newcommand{\vsat}{v_{\mathrm{sat}}}
\newcommand{\Voc}{V_{oc}}
\newcommand{\Isc}{I_{sc}}
\newcommand{\diapo}[1]{{\small\textsf{[diapo~#1]}}}
\newcommand{\dd}{\,\mathrm{d}}

% Distintivo para resultados que ya están en la hoja de fórmulas oficial.
% Uso: en el título de una caja, \begin{clave}[Título \enhoja], o en texto con \enhojatexto.
\newcommand{\enhojatexto}{{\scriptsize\sffamily\colorbox{verdeclaro}{\textcolor{verde!70!black}{\ensuremath{\checkmark}\,en la hoja}}}}
\newcommand{\enhoja}{\hfill{\scriptsize\sffamily\colorbox{white}{\textcolor{verde!70!black}{\ensuremath{\checkmark}\,en la hoja}}}}

% =====================================================================
%  Constantes oficiales de la materia ("Tabla de constantes.pdf").
%  Usar SIEMPRE estos valores en los ejemplos numéricos, para que
%  coincidan con los que se usan en clase y en el examen.
% =====================================================================

% ---------- Constantes físicas ----------
\newcommand{\qval}{\num{1.60e-19}}          % C
\newcommand{\kval}{\num{8.61e-5}}           % eV/K
\newcommand{\kJval}{\num{1.38e-23}}         % J/K
\newcommand{\hval}{\num{6.62e-34}}          % J s
\newcommand{\heVval}{\num{4.13e-15}}        % eV s
\newcommand{\epsval}{\num{8.85e-14}}        % F/cm
\newcommand{\VTval}{\SI{26}{mV}}            % kT/q a 300 K (0.0259 V)
\newcommand{\epsrox}{3.9}                   % SiO2

% ---------- Silicio ----------
\newcommand{\epsrSi}{11.8}
\newcommand{\niSi}{\num{1.0e10}}            % cm^-3
\newcommand{\munSi}{1350}                   % cm^2/(V s)
\newcommand{\mupSi}{480}                    % cm^2/(V s)
\newcommand{\EGSi}{1.12}                    % eV
\newcommand{\NCSi}{\num{3.22e19}}           % cm^-3
\newcommand{\NVSi}{\num{1.83e19}}           % cm^-3

% ---------- Germanio ----------
\newcommand{\epsrGe}{15.8}
\newcommand{\niGe}{\num{2.4e13}}
\newcommand{\munGe}{3900}
\newcommand{\mupGe}{1900}
\newcommand{\EGGe}{0.66}

% ---------- GaAs ----------
\newcommand{\epsrGaAs}{13.1}
\newcommand{\niGaAs}{\num{2.1e6}}
\newcommand{\munGaAs}{8800}
\newcommand{\mupGaAs}{400}
\newcommand{\EGGaAs}{1.42}

% ---------- Tabla para incluir en los apuntes ----------
\newcommand{\tablaconstantes}{%
\begin{tabular}{lccc}
\toprule
 & Si & Ge & GaAs \\
\midrule
$\epsilon_r$                     & 11.8 & 15.8 & 13.1 \\
$\mu_n$ (\si{cm^2/V.s})          & 1350 & 3900 & 8800 \\
$\mu_p$ (\si{cm^2/V.s})          & 480  & 1900 & 400 \\
$n_i$ (\si{cm^{-3}}, 300\,K)     & \num{1.0e10} & \num{2.4e13} & \num{2.1e6} \\
$E_G$ (\si{\eV})                 & 1.12 & 0.66 & 1.42 \\
\bottomrule
\end{tabular}}


% ---------- Estilos TikZ comunes ----------
\tikzset{
  banda/.style={line width=1.1pt},
  bandac/.style={banda,colEc},
  bandav/.style={banda,colEv},
  fermi/.style={dashed,gris,line width=0.6pt},
  elec/.style={circle,fill=colElec,inner sep=0pt,minimum size=5pt},
  hueco/.style={circle,draw=colHueco,fill=white,line width=0.9pt,inner sep=0pt,minimum size=5pt},
  foton/.style={colFoton,line width=0.9pt,decorate,
    decoration={snake,amplitude=1.6pt,segment length=5pt,post length=3pt},-{Stealth[length=5pt]}},
  flecha/.style={-{Stealth[length=5pt]},line width=0.7pt},
}

% =====================================================================

% ---------- Idioma y tipografía ----------
\usepackage[T1]{fontenc}
\usepackage[utf8]{inputenc}
\usepackage[spanish,es-nodecimaldot,es-noshorthands,es-tabla]{babel}
\usepackage{lmodern}
\usepackage{microtype}
\usepackage[a4paper,margin=2.3cm,headheight=14pt]{geometry}
\usepackage{parskip}
\setlength{\parskip}{0.55em}

% ---------- Matemática y unidades ----------
\usepackage{amsmath,amssymb,mathtools}
\usepackage{bm}
\usepackage[per-mode=symbol,output-decimal-marker={.},separate-uncertainty]{siunitx}
\sisetup{range-phrase={\text{--}},range-units=single}
\DeclareSIUnit{\eV}{eV}

% ---------- Gráficos ----------
\usepackage{graphicx}
\usepackage{xcolor}
\usepackage{tikz}
\usetikzlibrary{arrows.meta,decorations.pathmorphing,decorations.markings,
  patterns,calc,positioning,shapes.geometric,fit,backgrounds,shadings}
\usepackage{pgfplots}
\pgfplotsset{compat=1.18}
\usepgfplotslibrary{fillbetween}
\usepackage[siunitx,RPvoltages]{circuitikz}

% ---------- Tablas, listas, figuras ----------
\usepackage{booktabs,array,tabularx,multirow}
\usepackage{enumitem}
\setlist{itemsep=0.2em,topsep=0.3em}
\usepackage{caption,subcaption}
\captionsetup{font=small,labelfont=bf,format=hang}
\usepackage{float}
\usepackage{wrapfig}
\usepackage{needspace}

% ---------- Colores ----------
\definecolor{azul}{HTML}{1F4E79}
\definecolor{azulclaro}{HTML}{DCE9F5}
\definecolor{verde}{HTML}{2E7D32}
\definecolor{verdeclaro}{HTML}{E3F2E1}
\definecolor{naranja}{HTML}{C55A11}
\definecolor{naranjaclaro}{HTML}{FBE9DB}
\definecolor{rojo}{HTML}{B71C1C}
\definecolor{rojoclaro}{HTML}{FDECEA}
\definecolor{violeta}{HTML}{5E35B1}
\definecolor{violetaclaro}{HTML}{EDE7F6}
\definecolor{gris}{HTML}{555555}
% colores para dibujos de bandas / portadores
\colorlet{colEc}{azul}
\colorlet{colEv}{rojo}
\colorlet{colElec}{azul}
\colorlet{colHueco}{rojo}
\colorlet{colFoton}{naranja}

% ---------- Cajas didácticas ----------
\usepackage[most]{tcolorbox}
\tcbset{
  base/.style={enhanced,breakable,arc=2pt,boxrule=0.6pt,left=6pt,right=6pt,
    top=4pt,bottom=4pt,fonttitle=\bfseries\small,before skip=8pt,after skip=8pt}
}
% Definición
\newtcolorbox{definicion}[1][]{base,colback=azulclaro,colframe=azul,
  title={Definición\ifx\relax#1\relax\else: #1\fi}}
% Resultado / fórmula clave
\newtcolorbox{clave}[1][]{base,colback=verdeclaro,colframe=verde,
  title={\ifx\relax#1\relax Resultado clave\else #1\fi}}
% Ejemplo resuelto
\newtcolorbox{ejemplo}[1][]{base,colback=white,colframe=violeta,
  coltitle=white,colbacktitle=violeta,
  title={Ejemplo\ifx\relax#1\relax\else: #1\fi}}
% Corrección / aclaración respecto de las diapositivas
\newtcolorbox{ojo}[1][]{base,colback=rojoclaro,colframe=rojo,
  title={\ensuremath{\triangle}\ Ojo\ifx\relax#1\relax\else: #1\fi}}
% Intuición física / comentario
\newtcolorbox{intuicion}[1][]{base,colback=naranjaclaro,colframe=naranja,
  title={Idea física\ifx\relax#1\relax\else: #1\fi}}

% ---------- Encabezados ----------
\usepackage{fancyhdr}
\pagestyle{fancy}
\fancyhf{}
\renewcommand{\headrulewidth}{0.4pt}
\fancyhead[L]{\small\color{gris}Física Electrónica}
\fancyhead[R]{\small\color{gris}\unidadtitulo}
\fancyfoot[C]{\small\thepage}
\newcommand{\unidadtitulo}{}
\newcommand{\unidad}[1]{\renewcommand{\unidadtitulo}{#1}}

\usepackage{titlesec}
\titleformat{\section}{\Large\bfseries\color{azul}}{\thesection}{0.6em}{}
\titleformat{\subsection}{\large\bfseries\color{azul}}{\thesubsection}{0.6em}{}
\titleformat{\subsubsection}{\normalsize\bfseries\color{azul!80!black}}{\thesubsubsection}{0.6em}{}

\usepackage[hidelinks,bookmarksnumbered]{hyperref}
\usepackage{bookmark}

% ---------- Macros físicas ----------
\newcommand{\EG}{E_G}
\newcommand{\Ec}{E_c}
\newcommand{\Ev}{E_v}
\newcommand{\EF}{E_F}
\newcommand{\Eph}{E_{\mathrm{ph}}}
\newcommand{\lG}{\lambda_G}
\newcommand{\VT}{V_T}
\newcommand{\Popt}{P_{\mathrm{opt}}}
\newcommand{\Resp}{\mathcal{R}}
\newcommand{\vsat}{v_{\mathrm{sat}}}
\newcommand{\Voc}{V_{oc}}
\newcommand{\Isc}{I_{sc}}
\newcommand{\diapo}[1]{{\small\textsf{[diapo~#1]}}}
\newcommand{\dd}{\,\mathrm{d}}

% Distintivo para resultados que ya están en la hoja de fórmulas oficial.
% Uso: en el título de una caja, \begin{clave}[Título \enhoja], o en texto con \enhojatexto.
\newcommand{\enhojatexto}{{\scriptsize\sffamily\colorbox{verdeclaro}{\textcolor{verde!70!black}{\ensuremath{\checkmark}\,en la hoja}}}}
\newcommand{\enhoja}{\hfill{\scriptsize\sffamily\colorbox{white}{\textcolor{verde!70!black}{\ensuremath{\checkmark}\,en la hoja}}}}

% =====================================================================
%  Constantes oficiales de la materia ("Tabla de constantes.pdf").
%  Usar SIEMPRE estos valores en los ejemplos numéricos, para que
%  coincidan con los que se usan en clase y en el examen.
% =====================================================================

% ---------- Constantes físicas ----------
\newcommand{\qval}{\num{1.60e-19}}          % C
\newcommand{\kval}{\num{8.61e-5}}           % eV/K
\newcommand{\kJval}{\num{1.38e-23}}         % J/K
\newcommand{\hval}{\num{6.62e-34}}          % J s
\newcommand{\heVval}{\num{4.13e-15}}        % eV s
\newcommand{\epsval}{\num{8.85e-14}}        % F/cm
\newcommand{\VTval}{\SI{26}{mV}}            % kT/q a 300 K (0.0259 V)
\newcommand{\epsrox}{3.9}                   % SiO2

% ---------- Silicio ----------
\newcommand{\epsrSi}{11.8}
\newcommand{\niSi}{\num{1.0e10}}            % cm^-3
\newcommand{\munSi}{1350}                   % cm^2/(V s)
\newcommand{\mupSi}{480}                    % cm^2/(V s)
\newcommand{\EGSi}{1.12}                    % eV
\newcommand{\NCSi}{\num{3.22e19}}           % cm^-3
\newcommand{\NVSi}{\num{1.83e19}}           % cm^-3

% ---------- Germanio ----------
\newcommand{\epsrGe}{15.8}
\newcommand{\niGe}{\num{2.4e13}}
\newcommand{\munGe}{3900}
\newcommand{\mupGe}{1900}
\newcommand{\EGGe}{0.66}

% ---------- GaAs ----------
\newcommand{\epsrGaAs}{13.1}
\newcommand{\niGaAs}{\num{2.1e6}}
\newcommand{\munGaAs}{8800}
\newcommand{\mupGaAs}{400}
\newcommand{\EGGaAs}{1.42}

% ---------- Tabla para incluir en los apuntes ----------
\newcommand{\tablaconstantes}{%
\begin{tabular}{lccc}
\toprule
 & Si & Ge & GaAs \\
\midrule
$\epsilon_r$                     & 11.8 & 15.8 & 13.1 \\
$\mu_n$ (\si{cm^2/V.s})          & 1350 & 3900 & 8800 \\
$\mu_p$ (\si{cm^2/V.s})          & 480  & 1900 & 400 \\
$n_i$ (\si{cm^{-3}}, 300\,K)     & \num{1.0e10} & \num{2.4e13} & \num{2.1e6} \\
$E_G$ (\si{\eV})                 & 1.12 & 0.66 & 1.42 \\
\bottomrule
\end{tabular}}


% ---------- Estilos TikZ comunes ----------
\tikzset{
  banda/.style={line width=1.1pt},
  bandac/.style={banda,colEc},
  bandav/.style={banda,colEv},
  fermi/.style={dashed,gris,line width=0.6pt},
  elec/.style={circle,fill=colElec,inner sep=0pt,minimum size=5pt},
  hueco/.style={circle,draw=colHueco,fill=white,line width=0.9pt,inner sep=0pt,minimum size=5pt},
  foton/.style={colFoton,line width=0.9pt,decorate,
    decoration={snake,amplitude=1.6pt,segment length=5pt,post length=3pt},-{Stealth[length=5pt]}},
  flecha/.style={-{Stealth[length=5pt]},line width=0.7pt},
}

\unidad{Unidades 7 y 8 -- Transistor bipolar (BJT)}

\newcommand{\IFO}{I_{F0}}
\newcommand{\IRO}{I_{R0}}
\newcommand{\aF}{\alpha_F}
\newcommand{\aR}{\alpha_R}
\newcommand{\aT}{\alpha_T}
\newcommand{\bdc}{\beta_{dc}}
\newcommand{\adc}{\alpha_{dc}}
\newcommand{\eVEB}{\left(e^{V_{EB}/V_T}-1\right)}
\newcommand{\eVCB}{\left(e^{V_{CB}/V_T}-1\right)}

\title{\color{azul}\textbf{Transistor bipolar de juntura (BJT)}\\[0.3em]
  \large Funcionamiento, Ebers--Moll, efectos secundarios, saturación y modelos de pequeña señal}
\author{Física Electrónica -- Resumen de las unidades 7 y 8}
\date{Basado en las clases (BJT 1.\textsuperscript{a} y 2.\textsuperscript{a} parte) y en Pierret,
\emph{Semiconductor Device Fundamentals}, caps.~10, 11 y 12}

\begin{document}
\maketitle
\thispagestyle{empty}

\begin{tcolorbox}[base,colback=white,colframe=gris,title={Cómo leer este apunte}]
\small
\begin{itemize}
  \item Cajas \textcolor{azul}{\textbf{azules}}: definiciones; \textcolor{verde}{\textbf{verdes}}: resultados
  clave; \textcolor{naranja}{\textbf{naranjas}}: intuición física; \textcolor{violeta}{\textbf{violetas}}: ejemplos;
  \textcolor{rojo}{\textbf{rojas}} (\emph{Ojo}): errores o imprecisiones de las diapositivas y confusiones típicas.
  \item \diapo{1.n} y \diapo{2.n} indican la diapositiva $n$ de la 1.\textsuperscript{a} o 2.\textsuperscript{a}
  presentación de BJT.
  \item El sello \enhojatexto{} indica que el resultado figura en la hoja de fórmulas oficial.
  \item \textbf{Se trabaja con un transistor PNP}, igual que la hoja de fórmulas y Pierret. Algunas
  diapositivas (tomadas de Sedra--Smith) usan NPN: la traducción está en la tabla~\ref{tab:pnp-npn}.
  \item Los ejemplos numéricos usan los valores de la \emph{tabla de constantes} de la materia.
\end{itemize}
\end{tcolorbox}

\tableofcontents
\newpage

% =====================================================================
\section{Fundamentos}
% =====================================================================

\subsection{Un poco de historia}

El primer transistor (de punta de contacto, sobre germanio) lo construyeron Bardeen y
Brattain en los Bell Labs en diciembre de 1947; Shockley propuso poco después el
transistor \emph{de juntura}, que es el que estudiamos. Los tres recibieron el Nobel de
Física en 1956 \diapo{1.2--3}. El transistor reemplazó a las válvulas de vacío: funciona con
menos tensión, es más chico, robusto y barato, y se puede integrar.

\subsection{Estructura, símbolo y convenciones}

Un BJT son \textbf{dos junturas PN muy próximas} que comparten una región central angosta,
la \textbf{base}. En un PNP: emisor p$^+$, base n, colector p. El dopaje es deliberadamente
asimétrico:
\[
  N_E \gg N_B > N_C
  \qquad(\text{típico: } \SI{e18}{} \text{ -- } \SI{e19}{},\ \SI{e16}{} \text{ -- } \SI{e17}{},\
  \SI{e15}{} \text{ -- } \SI{e16}{cm^{-3}}).
\]
El emisor ``emite'' (inyecta) portadores en la base y el colector los ``colecta''. Por eso
el dispositivo \emph{no es simétrico}: intercambiar emisor y colector da un transistor mucho
peor (modo inverso).

\begin{figure}[H]
\centering
\begin{minipage}{0.52\textwidth}
\centering
\begin{tikzpicture}[scale=0.95]
  \fill[rojoclaro] (0,0) rectangle (2.4,1.4);
  \fill[azulclaro] (2.4,0) rectangle (3.4,1.4);
  \fill[rojo!8] (3.4,0) rectangle (6.4,1.4);
  \draw (0,0) rectangle (6.4,1.4) (2.4,0) -- (2.4,1.4) (3.4,0) -- (3.4,1.4);
  \node at (1.2,0.95) {p$^+$}; \node[font=\scriptsize] at (1.2,0.45) {emisor};
  \node at (2.9,0.95) {n}; \node[font=\scriptsize] at (2.9,0.45) {base};
  \node at (4.9,0.95) {p}; \node[font=\scriptsize] at (4.9,0.45) {colector};
  \draw (0,0.7) -- (-0.6,0.7) node[left] {E};
  \draw (6.4,0.7) -- (7.0,0.7) node[right] {C};
  \draw (2.9,0) -- (2.9,-0.6) node[below] {B};
  \draw[<->] (2.4,1.65) -- node[above,font=\small] {$W_B$} (3.4,1.65);
  \node[font=\scriptsize,anchor=east] at (2.35,-0.35) {J$_{E}$ (E--B)};
  \node[font=\scriptsize,anchor=west] at (3.45,-0.35) {J$_{C}$ (C--B)};
\end{tikzpicture}
\end{minipage}%
\begin{minipage}{0.46\textwidth}
\centering
\begin{circuitikz}[scale=0.9,transform shape]
  \draw (0,0) node[pnp](Q){};
  \draw (Q.E) -- ++(0,0.4) node[above] {E};
  \draw (Q.C) -- ++(0,-0.4) node[below] {C};
  \draw (Q.B) -- ++(-0.4,0) node[left] {B};
  \draw[flecha,azul] ($(Q.E)+(0.3,0.5)$) -- node[right,font=\small] {$I_E$} ($(Q.E)+(0.3,0.05)$);
  \draw[flecha,azul] ($(Q.C)+(0.3,-0.05)$) -- node[right,font=\small] {$I_C$} ($(Q.C)+(0.3,-0.5)$);
  \draw[flecha,azul] ($(Q.B)+(-0.2,0.25)$) -- node[above,font=\small] {$I_B$} ($(Q.B)+(-0.8,0.25)$);
  \begin{scope}[xshift=3.2cm]
  \draw (0,0) node[npn](Q2){};
  \draw (Q2.C) -- ++(0,0.4) node[above] {C};
  \draw (Q2.E) -- ++(0,-0.4) node[below] {E};
  \draw (Q2.B) -- ++(-0.4,0) node[left] {B};
  \node[font=\small\bfseries] at (0,-2.0) {NPN};
  \end{scope}
  \node[font=\small\bfseries] at (0,-2.0) {PNP};
\end{circuitikz}
\end{minipage}
\caption{Estructura de un PNP (no está a escala: la base es mucho más angosta) y símbolos
\diapo{1.4}. La flecha del símbolo está en el emisor y apunta en el sentido de la corriente
convencional por la juntura E--B en directa: hacia adentro en el PNP, hacia afuera en el NPN.}
\label{fig:estructura}
\end{figure}

\begin{clave}[Convenciones (PNP)]
\[
  I_E = I_C + I_B,
  \qquad
  V_{EC} = V_{EB} + V_{BC} = V_{EB} - V_{CB}.
\]
$I_E$ entra por el emisor; $I_C$ e $I_B$ salen por colector y base. En funcionamiento normal
todas son positivas. Las tensiones $V_{EB}$ y $V_{CB}$ se miden respecto de la base.
\end{clave}

\paragraph{Configuraciones \diapo{1.5}.} El transistor tiene tres terminales, pero un circuito
amplificador tiene una malla de entrada y una de salida (cuatro terminales): uno de los tres
terminales es \textbf{común} a ambas. Hay entonces tres configuraciones: \emph{base común}
(entrada E, salida C), \emph{emisor común} (entrada B, salida C, la más usada) y
\emph{colector común} (entrada B, salida E; seguidor de emisor).

\subsection{Regiones de operación}

Cada juntura puede estar en directa o en inversa, lo que da cuatro regiones \diapo{1.6}:

\begin{table}[H]
\centering\small
\caption{Regiones de operación del BJT (PNP: ``directa'' en E--B significa $V_{EB}>0$; en
C--B, $V_{CB}>0$).}
\label{tab:regiones}
\begin{tabular}{lccl}
\toprule
Región & Juntura E--B & Juntura C--B & Uso \\
\midrule
\textbf{Activa directa (MAD)} & directa & inversa & amplificación \\
Saturación                   & directa & directa & llave cerrada \\
Corte                        & inversa & inversa & llave abierta \\
Activa inversa (MAI)         & inversa & directa & casi no se usa (ganancia muy baja) \\
\bottomrule
\end{tabular}
\end{table}

\subsection{¿Es lo mismo que dos diodos enfrentados?}

\textbf{No} \diapo{1.7}. Con dos diodos discretos unidos por un cable, la corriente del diodo
E--B (en directa) se recombina en el cable y en las zonas neutras: al diodo C--B (en
inversa) sólo le llega su pequeña corriente de saturación. El efecto transistor aparece
porque la base es \textbf{mucho más angosta que la longitud de difusión} de sus
minoritarios ($W\ll L_B$): los huecos que el emisor inyecta en la base la cruzan por
difusión \emph{casi sin recombinarse} y llegan a la juntura C--B, donde el campo los barre al
colector. Así, una juntura en inversa conduce una corriente grande, controlada por la
tensión de la otra juntura.

\subsection{Diagramas de bandas}

\begin{figure}[H]
\centering
\begin{tikzpicture}[x=0.95cm,y=1cm]
  % ---------- equilibrio ----------
  \node[font=\small\bfseries,anchor=west] at (-0.3,3.6) {(a) Equilibrio};
  \draw[bandac] (0,2.0) -- (2,2.0) to[out=0,in=180] (3,3.0) -- (4.2,3.0) to[out=0,in=180] (5.2,2.0) -- (7.5,2.0) node[right,font=\small] {$\Ec$};
  \draw[bandav] (0,0.9) -- (2,0.9) to[out=0,in=180] (3,1.9) -- (4.2,1.9) to[out=0,in=180] (5.2,0.9) -- (7.5,0.9) node[right,font=\small] {$\Ev$};
  \draw[fermi] (0,1.0) -- (7.5,1.0);
  \node[font=\scriptsize,anchor=east] at (0,1.05) {$\EF$};
  \foreach \x in {0.3,0.7,1.1,1.5} \node[hueco] at (\x,0.85) {};
  \foreach \x in {5.6,6.2,6.8} \node[hueco] at (\x,0.85) {};
  \foreach \x in {3.2,3.6,4.0} \node[elec] at (\x,3.05) {};
  \node[font=\scriptsize] at (1,0.3) {E (p$^+$)}; \node[font=\scriptsize] at (3.6,0.3) {B (n)};
  \node[font=\scriptsize] at (6.3,0.3) {C (p)};
  % ---------- MAD ----------
  \begin{scope}[xshift=8.4cm]
  \node[font=\small\bfseries,anchor=west] at (-0.3,3.6) {(b) Activa directa};
  \draw[bandac] (0,2.0) -- (2,2.0) to[out=0,in=180] (2.8,2.55) -- (4.2,2.55) to[out=0,in=180] (5.6,0.6) -- (7.5,0.6);
  \draw[bandav] (0,0.9) -- (2,0.9) to[out=0,in=180] (2.8,1.45) -- (4.2,1.45) to[out=0,in=180] (5.6,-0.5) -- (7.5,-0.5);
  \draw[fermi] (0,1.0) -- (2,1.0);
  \draw[fermi] (5.6,-0.4) -- (7.5,-0.4);
  \foreach \x in {0.3,0.7,1.1,1.5} \node[hueco] at (\x,0.85) {};
  \draw[flecha,colHueco,line width=0.9pt] (1.4,1.05) to[out=20,in=180] (2.8,1.3) -- (4.0,1.3) to[out=0,in=150] (5.3,-0.2);
  \node[colHueco,font=\scriptsize,align=center] at (3.5,0.65) {huecos inyectados\\cruzan la base};
  \node[font=\scriptsize] at (1,0.3) {E}; \node[font=\scriptsize] at (3.5,3.0) {B};
  \node[font=\scriptsize] at (6.6,-1.0) {C};
  \node[font=\scriptsize,align=center,gris] at (1.6,-0.4) {barrera E--B\\baja $qV_{EB}$};
  \node[font=\scriptsize,align=center,gris] at (6.6,1.6) {barrera C--B\\sube $q|V_{CB}|$};
  \end{scope}
\end{tikzpicture}
\caption{Diagrama de bandas de un PNP \diapo{1.8--1.10}. Los huecos tienen más energía cuanto
más \emph{abajo} están en el diagrama: en MAD ``caen hacia arriba'' desde el emisor hacia la
base (barrera reducida) y luego ``ruedan'' por la pendiente de la juntura C--B hacia el
colector. En la base no hay campo (es neutra): los huecos la cruzan por \textbf{difusión}.}
\label{fig:bandas}
\end{figure}

\subsection{Flujos de portadores en activa directa}

\begin{figure}[H]
\centering
\begin{tikzpicture}[x=1cm,y=1cm]
  \fill[rojoclaro] (0,0) rectangle (3.2,3.2);
  \fill[gris!20] (3.2,0) rectangle (3.6,3.2);
  \fill[azulclaro] (3.6,0) rectangle (5.6,3.2);
  \fill[gris!20] (5.6,0) rectangle (6.4,3.2);
  \fill[rojo!8] (6.4,0) rectangle (9.6,3.2);
  \draw (0,0) rectangle (9.6,3.2);
  \foreach \x in {3.2,3.6,5.6,6.4} \draw[densely dashed] (\x,0) -- (\x,3.2);
  \node[font=\small] at (1.2,3.45) {E (p$^+$)}; \node[font=\small] at (4.6,3.45) {B (n)};
  \node[font=\small] at (8.4,3.45) {C (p)};
  \draw[<->] (3.6,-0.25) -- node[below,font=\small,pos=0.3] {$W$} (5.6,-0.25);
  \draw[<->] (3.2,-0.85) -- node[below,font=\small,pos=0.25] {$W_B$} (6.4,-0.85);
  % huecos emitidos -> colector (I_Ep -> I_Cp)
  \foreach \y in {2.8,2.5,2.2} \draw[flecha,colHueco,line width=0.9pt] (1.6,\y) -- (8.2,\y);
  \node[colHueco,font=\small,anchor=east] at (1.55,2.5) {$I_{Ep}$};
  \node[colHueco,font=\small,anchor=west] at (8.25,2.5) {$I_{Cp}$};
  % hueco que se recombina en la base
  \draw[flecha,colHueco,line width=0.9pt] (1.6,1.75) -- (4.6,1.75);
  \node[hueco] at (4.75,1.75) {}; \node[elec] at (4.75,1.45) {};
  \draw[flecha,colElec] (4.75,-0.55) -- (4.75,1.3);
  \node[colElec,font=\scriptsize,anchor=west,fill=azulclaro,inner sep=1pt] at (4.8,1.2) {$I_{B2}$};
  % electrones base -> emisor (I_En)
  \draw[flecha,colElec,line width=0.9pt] (4.2,1.05) -- (0.3,1.05);
  \node[colElec,font=\small,anchor=south west] at (0.3,1.05) {$I_{En}=I_{B1}$};
  % electrones colector -> base (I_Cn)
  \draw[flecha,colElec,line width=0.9pt] (9.3,0.45) -- (5.2,0.45);
  \node[colElec,font=\small,anchor=south east] at (9.3,0.45) {$I_{Cn}=I_{B3}$};
  % terminales
  \draw[line width=1.5pt] (0,0) -- (0,3.2) (9.6,0) -- (9.6,3.2);
  \draw (4.75,0) -- (4.75,-1.9) node[below] {B};
  \draw[flecha] (5.05,-1.4) -- node[right,font=\small] {$I_B$} (5.05,-1.85);
  \draw (0,1.6) -- (-0.7,1.6) node[left] {E};
  \draw[flecha] (-0.7,1.9) -- node[above,font=\small] {$I_E$} (-0.1,1.9);
  \draw (9.6,1.6) -- (10.3,1.6) node[right] {C};
  \draw[flecha] (9.7,1.9) -- node[above,font=\small] {$I_C$} (10.3,1.9);
\end{tikzpicture}
\caption{Componentes de corriente en un PNP en activa directa (Pierret, fig.~10.6 y
\diapo{1.11, 1.17}). Las flechas indican el movimiento de los portadores; la corriente de
electrones tiene sentido opuesto a su movimiento.}
\label{fig:flujos}
\end{figure}

\begin{itemize}
  \item $I_{Ep}$: huecos inyectados del emisor a la base. \textbf{Es la componente útil.}
  \item $I_{Cp}$: la parte de $I_{Ep}$ que llega al colector ($I_{Cp}\lesssim I_{Ep}$).
  \item $I_{B2}=I_{Ep}-I_{Cp}=I_{RB}$: huecos que se recombinan en la base; la base tiene que
  aportar los electrones para esa recombinación.
  \item $I_{En}=I_{B1}$: electrones inyectados de la base al emisor. Es inevitable (la juntura
  E--B en directa inyecta en ambos sentidos) pero \emph{no sirve}: se busca minimizarla con
  $N_E\gg N_B$.
  \item $I_{Cn}=I_{B3}$: corriente inversa de la juntura C--B (electrones térmicos del
  colector que pasan a la base). Muy chica.
\end{itemize}

\begin{clave}[Balance de corrientes]
\[
  I_E = I_{Ep}+I_{En},\qquad I_C = I_{Cp}+I_{Cn},\qquad
  I_B = I_{En} + (I_{Ep}-I_{Cp}) - I_{Cn} = I_{B1}+I_{B2}-I_{B3}.
\]
\end{clave}

\begin{intuicion}[una corriente chica controla una grande]
En un PNP, la corriente de base está formada por \emph{electrones} (los que van al emisor y
los que se recombinan), mientras que la corriente principal emisor--colector son
\emph{huecos} \diapo{1.18}. Con un buen diseño $I_B$ es 100 a 300 veces menor que $I_C$.
Ojo con la interpretación: lo que controla físicamente a $I_C$ es la \textbf{tensión}
$V_{EB}$, que fija la altura de la barrera y por lo tanto cuántos huecos se inyectan
($I_C\propto e^{V_{EB}/V_T}$). Que $I_C=\beta I_B$ es una consecuencia: $I_B$ y $I_C$ crecen
juntas con $V_{EB}$, en una proporción fija.
\end{intuicion}

\begin{ojo}[``el BJT es un dispositivo controlado por corriente'' \diapo{1.31, 2.34}]
Es una descripción \emph{de circuito} útil (modelo con $\beta$), no la física del dispositivo.
La corriente de colector depende exponencialmente de $V_{EB}$ y \emph{no} de $I_B$ en sí
misma; por eso la transconductancia $g_m=I_C/V_T$ es tan grande y el modelo de pequeña señal
usa una fuente controlada por \emph{tensión} ($g_mv_{eb}$). La diferencia práctica con el FET
es que el BJT \emph{necesita} una corriente de entrada ($I_B\neq0$) para funcionar.
\end{ojo}

% =====================================================================
\section{Análisis formal del transistor ideal}
\label{sec:formal}
% =====================================================================

El objetivo es obtener las corrientes $I_E$, $I_C$, $I_B$ en función de $V_{EB}$ y $V_{CB}$.
El método es el mismo que para el diodo ideal: resolver la ecuación de difusión de
minoritarios en cada región cuasineutra y evaluar las corrientes de difusión en los bordes
de las zonas de carga espacial.

\subsection{Hipótesis y notación}

\begin{itemize}
  \item Unidimensional, estado estacionario, dopajes uniformes en cada región.
  \item Baja inyección en todas las regiones; sin generación ni recombinación en las ZCE.
  \item Sin campo en las regiones cuasineutras: los minoritarios se mueven sólo por difusión.
  \item Emisor y colector ``largos'' (mucho más anchos que $L_E$, $L_C$).
\end{itemize}

\begin{figure}[H]
\centering
\begin{tikzpicture}[x=1cm,y=1cm]
  \fill[rojoclaro] (0,0) rectangle (4,1.2);
  \fill[gris!20] (4,0) rectangle (4.5,1.2);
  \fill[azulclaro] (4.5,0) rectangle (7,1.2);
  \fill[gris!20] (7,0) rectangle (7.8,1.2);
  \fill[rojo!8] (7.8,0) rectangle (11.8,1.2);
  \draw (0,0) rectangle (11.8,1.2);
  \node at (2,0.6) {Emisor p$^+$}; \node at (5.75,0.6) {Base n}; \node at (9.8,0.6) {Colector p};
  \draw[flecha] (4,-0.4) -- (1.2,-0.4) node[left] {$x''$};
  \node[below] at (4,-0.45) {$0$};
  \draw[flecha] (4.5,-1.1) -- (7.5,-1.1) node[right] {$x$};
  \node[below] at (4.5,-1.15) {$0$}; \node[below] at (7,-1.15) {$W$};
  \draw (7,-1.0) -- (7,-1.2);
  \draw[flecha] (7.8,-0.4) -- (10.6,-0.4) node[right] {$x'$};
  \node[below] at (7.8,-0.45) {$0$};
  \draw[<->] (4,1.45) -- node[above,font=\small] {$W_B$ (metalúrgico)} (7.8,1.45);
  \node[font=\scriptsize,align=left,anchor=north west] at (0,-1.7) {%
    $N_E=N_{AE}$, $D_E=D_N$, $\tau_E=\tau_n$, $L_E=L_N$, $n_{E0}=n_i^2/N_E$};
  \node[font=\scriptsize,align=left,anchor=north west] at (0,-2.15) {%
    $N_B=N_{DB}$, $D_B=D_P$, $\tau_B=\tau_p$, $L_B=L_P$, $p_{B0}=n_i^2/N_B$};
  \node[font=\scriptsize,align=left,anchor=north west] at (0,-2.6) {%
    $N_C=N_{AC}$, $D_C=D_N$, $\tau_C=\tau_n$, $L_C=L_N$, $n_{C0}=n_i^2/N_C$};
\end{tikzpicture}
\caption{Coordenadas y notación \diapo{1.12} (Pierret, fig.~11.2). Cada región tiene su propio
eje con origen en el borde de la ZCE; $x''$ crece hacia el contacto de emisor y $x'$ hacia el de
colector. $W$ es el ancho \emph{cuasineutro} de la base.}
\label{fig:coordenadas}
\end{figure}

\begin{ojo}[$W\neq W_B$ \diapo{1.12, 1.16}]
\[
  W = W_B - x_{nEB} - x_{nCB}
\]
$W_B$ es el ancho metalúrgico (entre junturas) y $W$ el ancho \emph{cuasineutro}, lo que queda
después de restar la parte de las dos ZCE que penetra en la base. $W$ \textbf{depende de las
tensiones aplicadas}: es el origen del efecto Early (sección~\ref{sec:early}).
\end{ojo}

\subsection{Emisor}

Ecuación de difusión (estacionaria, sin generación) y condiciones de contorno:
\[
  0 = D_E\frac{\dd^2\Delta n_E}{\dd x''^2} - \frac{\Delta n_E}{\tau_E},
  \qquad
  \Delta n_E(0) = n_{E0}\eVEB,
  \qquad
  \Delta n_E(x''\to\infty)=0 .
\]
La primera condición es la \textbf{ley de la juntura} aplicada a E--B. La solución es la del
diodo largo:
\[
  \Delta n_E(x'') = n_{E0}\eVEB e^{-x''/L_E},
  \qquad L_E=\sqrt{D_E\tau_E}.
\]
La corriente de electrones (que fluye en el sentido $+x$, de E hacia B, igual que $I_E$) es
\[
  I_{En} = -qAD_E\left.\frac{\dd\Delta n_E}{\dd x''}\right|_{x''=0}
         = qA\frac{D_E}{L_E}\,n_{E0}\eVEB .
\]

\subsection{Colector}

Análogamente, con la ley de la juntura en C--B:
\[
  \Delta n_C(x') = n_{C0}\eVCB e^{-x'/L_C},
\]
\[
  I_{Cn} = qAD_C\left.\frac{\dd\Delta n_C}{\dd x'}\right|_{x'=0}
         = -qA\frac{D_C}{L_C}\,n_{C0}\eVCB .
\]
En MAD $V_{CB}<0$, el paréntesis vale $\approx-1$ e $I_{Cn}\approx+qA(D_C/L_C)n_{C0}$: una
corriente de saturación inversa, chica y positiva.

\begin{ojo}[signo de $I_{Cn}$ \diapo{1.14}]
La diapositiva escribe $I_{Cn}=-qAD_C\,\dd\Delta n_C/\dd x'$, pero el resultado final que
muestra ($-qA\frac{D_C}{L_C}n_{C0}(\cdots)$, que coincide con la hoja) corresponde a
$+qAD_C\,\dd\Delta n_C/\dd x'$. La corriente de electrones en la dirección $+x'$ es
$J_n=+qD_n\,\dd n/\dd x'$; el signo menos aparece en el emisor porque allí $x''$ apunta al revés.
\end{ojo}

\subsection{Base}

\[
  0 = D_B\frac{\dd^2\Delta p_B}{\dd x^2} - \frac{\Delta p_B}{\tau_B},
  \qquad
  \Delta p_B(0) = p_{B0}\eVEB,
  \qquad
  \Delta p_B(W) = p_{B0}\eVCB .
\]
Ahora la región es \emph{corta} (ancho finito $W$) y hay una condición en cada borde. La
solución general $Ae^{x/L_B}+Be^{-x/L_B}$, escrita con funciones hiperbólicas
($\sinh x=\frac{e^x-e^{-x}}{2}$, $\cosh x=\frac{e^x+e^{-x}}{2}$), queda
\begin{equation}
  \Delta p_B(x) = \Delta p_B(0)\,\frac{\sinh\left[(W-x)/L_B\right]}{\sinh(W/L_B)}
                + \Delta p_B(W)\,\frac{\sinh(x/L_B)}{\sinh(W/L_B)} .
  \label{eq:pB-exacta}
\end{equation}
(Es fácil verificar las condiciones de contorno: en $x=0$ el segundo término se anula y el
primero vale $\Delta p_B(0)$; en $x=W$, al revés.)

Las corrientes de huecos en los bordes son corrientes de difusión, $I_p=-qAD_B\,\dd\Delta p_B/\dd x$:

\begin{clave}[Corrientes de huecos en la base \enhoja]
\[
  I_{Ep} = qA\frac{D_B}{L_B}p_{B0}\left[\frac{e^{V_{EB}/V_T}-1}{\tanh(W/L_B)}
         - \frac{e^{V_{CB}/V_T}-1}{\sinh(W/L_B)}\right]
\]
\[
  I_{Cp} = qA\frac{D_B}{L_B}p_{B0}\left[\frac{e^{V_{EB}/V_T}-1}{\sinh(W/L_B)}
         - \frac{e^{V_{CB}/V_T}-1}{\tanh(W/L_B)}\right]
\]
\end{clave}

\begin{ojo}[signo de las corrientes de huecos \diapo{1.15--1.16}]
La diapositiva 15 dice ``condiciones de contorno en el colector'' para las de la base. Y en
la 16 se escribe $I_{Ep}=+qAD_B\,\dd\Delta p_B/\dd x|_{x=0}$: falta el signo menos de la
corriente de difusión ($J_p=-qD_p\,\dd p/\dd x$). Como $\Delta p_B$ \emph{decrece} con $x$, la
derivada es negativa y la corriente es positiva. Los resultados finales de la diapo son
correctos.
\end{ojo}

\subsection{Corrientes en los terminales}

\begin{clave}[Corrientes del BJT ideal \enhoja]
\[
  I_E = I_{Ep}+I_{En},\qquad I_C = I_{Cp}+I_{Cn},\qquad
  I_B = I_E - I_C = I_{En}+I_{RB}-I_{Cn},\quad I_{RB}=I_{Ep}-I_{Cp}.
\]
\end{clave}

\subsection{Aproximación de base angosta}
\label{sec:base-angosta}

En un transistor real $W\ll L_B$ (típicamente $W\sim\SI{1}{\micro m}$ y $L_B\sim$ decenas
de \si{\micro m}). Con $\sinh u\approx u$, $\cosh u\approx 1+u^2/2$ para $u\ll1$,
\eqref{eq:pB-exacta} se vuelve una \textbf{recta}:
\begin{equation}
  \Delta p_B(x) \simeq \Delta p_B(0)\left(1-\frac{x}{W}\right) + \Delta p_B(W)\,\frac{x}{W}
  \qquad\text{\enhojatexto}
  \label{eq:pB-lineal}
\end{equation}
Físicamente: si casi no hay recombinación, la corriente de difusión es la misma en toda la
base, y eso exige un gradiente constante.

\begin{figure}[H]
\centering
\begin{tikzpicture}
\begin{axis}[
  width=0.9\textwidth,height=6.2cm,
  xmin=-3.2,xmax=4.4,ymin=-0.08,ymax=1.15,
  axis lines=none,clip=false,
]
% zonas
\fill[rojoclaro] (axis cs:-3.2,-0.08) rectangle (axis cs:-0.3,1.15);
\fill[gris!20] (axis cs:-0.3,-0.08) rectangle (axis cs:0,1.15);
\fill[azulclaro] (axis cs:0,-0.08) rectangle (axis cs:1,1.15);
\fill[gris!20] (axis cs:1,-0.08) rectangle (axis cs:1.5,1.15);
\fill[rojo!8] (axis cs:1.5,-0.08) rectangle (axis cs:4.4,1.15);
\draw (axis cs:-3.2,0) -- (axis cs:4.4,0);
% base: exacta (con W/L=1 exagerado) y lineal
\addplot[colHueco,line width=1.3pt,domain=0:1] {1-x};
\addplot[colHueco,dashed,line width=1pt,domain=0:1,samples=40] {sinh(1.6*(1-x))/sinh(1.6)};
% emisor: electrones (escala distinta)
\addplot[colElec,line width=1.3pt,domain=-3.2:-0.3,samples=60] {0.35*exp((x+0.3)/0.8)};
% colector: deficit de electrones (negativo pequeño)
\addplot[colElec,line width=1.3pt,domain=1.5:4.4,samples=60] {-0.06*exp(-(x-1.5)/0.7)};
\node[colHueco,font=\scriptsize,anchor=south west] at (axis cs:0.12,0.82) {$\Delta p_B(x)$};
\node[colHueco,font=\scriptsize,align=left,anchor=west] at (axis cs:0.55,0.62) {recta si $W\ll L_B$\\(punteada: $W\sim L_B$)};
\node[colElec,font=\scriptsize,anchor=south] at (axis cs:-1.6,0.16) {$\Delta n_E(x'')$};
\node[colElec,font=\scriptsize,anchor=north] at (axis cs:2.9,-0.03) {$\Delta n_C(x')\approx-n_{C0}$};
\node[font=\scriptsize,anchor=east] at (axis cs:0,1.0) {$\Delta p_B(0)$};
\node[font=\small] at (axis cs:-1.75,1.07) {E}; \node[font=\small] at (axis cs:0.5,1.07) {B};
\node[font=\small] at (axis cs:2.95,1.07) {C};
\node[font=\scriptsize,anchor=north] at (axis cs:0,-0.08) {$0$};
\node[font=\scriptsize,anchor=north] at (axis cs:1,-0.08) {$W$};
\end{axis}
\end{tikzpicture}
\caption{Excesos de minoritarios en MAD (Pierret, fig.~11.3; \diapo{2.1}). Las escalas de
cada región son distintas: como $N_E\gg N_B$, $n_{E0}\ll p_{B0}$ y la inyección hacia el emisor
es mucho menor. En el colector (inversa) hay un pequeño \emph{defecto} de electrones.}
\label{fig:perfiles}
\end{figure}

En MAD ($V_{CB}\ll -V_T$, $\Delta p_B(W)\approx-p_{B0}\approx0$) la corriente de colector
queda
\[
  I_C \simeq I_{Cp} \simeq qA\frac{D_B}{W}\,p_{B0}\,e^{V_{EB}/V_T}
     = \underbrace{\frac{qAD_Bn_i^2}{N_BW}}_{I_S}\,e^{V_{EB}/V_T}.
\]
Dos conclusiones importantes:
\begin{enumerate}
  \item $I_C$ depende exponencialmente de $V_{EB}$ y es \textbf{independiente de $V_{CB}$}
  (mientras C--B esté en inversa) \diapo{1.10}.
  \item $I_C\propto 1/W$: si $W$ cambia con la tensión, cambia $I_C$ (efecto Early).
\end{enumerate}

\paragraph{Diodo corto \diapo{1.30}.} Si además el emisor y el colector son más cortos que
sus longitudes de difusión (lo habitual en transistores integrados), se puede despreciar la
recombinación en las tres regiones: todos los perfiles son rectas y en las expresiones
$L_E\to W_E'$, $L_C\to W_C'$ (anchos cuasineutros de emisor y colector), por ejemplo
\[
  I_{F0} = qAn_i^2\left(\frac{D_E}{N_EW_E'}+\frac{D_B}{N_BW}\right).
\]
Los anchos son parámetros de diseño mucho más controlables que las longitudes de difusión.

% =====================================================================
\section{Parámetros de rendimiento}
\label{sec:parametros}
% =====================================================================

La corriente de colector ``útil'' es la de huecos que el emisor inyecta y que llega al
colector. Entre la corriente total de emisor y esa corriente hay dos pérdidas, que se miden
con dos factores \diapo{1.24}:

\begin{definicion}[eficiencia de emisión y factor de transporte \enhoja]
\[
  \gamma = \frac{I_{Ep}}{I_E} = \frac{I_{Ep}}{I_{Ep}+I_{En}}
  \qquad\text{(qué fracción de $I_E$ son huecos inyectados en la base)}
\]
\[
  \aT = \frac{I_{Cp}}{I_{Ep}}
  \qquad\text{(qué fracción de esos huecos cruza la base sin recombinarse)}
\]
\end{definicion}

\begin{definicion}[ganancias de corriente]
\[
  \adc = \frac{I_{Cp}}{I_E} = \gamma\,\aT \approx \frac{I_C}{I_E}
  \quad\text{(base común)},
  \qquad
  \bdc = \frac{I_C}{I_B} = \frac{\adc}{1-\adc}
  \quad\text{(emisor común)}.
\]
Como $\adc$ es apenas menor que 1, $\bdc$ es grande: $\adc=0.99\Rightarrow\bdc=99$;
$\adc=0.995\Rightarrow\bdc=199$. Pequeñas mejoras en $\adc$ producen grandes cambios en $\bdc$.
\end{definicion}

\subsection{Expresiones exactas}

Evaluando las corrientes de la sección~\ref{sec:formal} en MAD
($e^{V_{EB}/V_T}\gg1$, $\Delta p_B(W)\approx0$) \diapo{1.25}:
\[
  \gamma = \frac{1}{1+\dfrac{D_E}{D_B}\dfrac{L_B}{L_E}\dfrac{N_B}{N_E}\tanh\!\left(\dfrac{W}{L_B}\right)},
  \qquad
  \aT = \frac{1}{\cosh(W/L_B)},
\]
\[
  \adc = \frac{1}{\cosh\!\left(\frac{W}{L_B}\right)
         + \frac{D_E}{D_B}\frac{L_B}{L_E}\frac{N_B}{N_E}\sinh\!\left(\frac{W}{L_B}\right)},
  \qquad
  \bdc = \frac{1}{\cosh\!\left(\frac{W}{L_B}\right)
         + \frac{D_E}{D_B}\frac{L_B}{L_E}\frac{N_B}{N_E}\sinh\!\left(\frac{W}{L_B}\right)-1}.
\]
(Se usó $n_{E0}/p_{B0}=N_B/N_E$.)

\subsection{Aproximación de base angosta}

Con $W\ll L_B$: $\tanh u\approx u$, $\cosh u\approx1+u^2/2$, y $1/(1+\epsilon)\approx1-\epsilon$
\diapo{1.26}:

\begin{clave}[Parámetros con base angosta \enhoja]
\[
  \gamma \simeq \frac{1}{1+\dfrac{D_E}{D_B}\dfrac{N_B}{N_E}\dfrac{W}{L_E}},
  \qquad
  \aT \simeq \frac{1}{1+\dfrac{1}{2}\left(\dfrac{W}{L_B}\right)^2},
\]
\[
  \adc \simeq \frac{1}{1+\dfrac{D_E}{D_B}\dfrac{N_B}{N_E}\dfrac{W}{L_E}
          +\dfrac{1}{2}\left(\dfrac{W}{L_B}\right)^2},
  \qquad
  \bdc \simeq \frac{1}{\dfrac{D_E}{D_B}\dfrac{N_B}{N_E}\dfrac{W}{L_E}
          +\dfrac{1}{2}\left(\dfrac{W}{L_B}\right)^2}.
\]
\end{clave}

\begin{intuicion}[cómo se diseña un buen BJT]
\begin{itemize}
  \item \textbf{$N_E\gg N_B$} (emisor mucho más dopado que la base) hace $\gamma\to1$: la
  juntura E--B inyecta casi sólo huecos hacia la base y casi no electrones hacia el emisor.
  \item \textbf{$W\ll L_B$} (base angosta, con buena vida media) hace $\aT\to1$: los huecos
  cruzan sin tiempo de recombinarse. Notar que $1-\aT\propto(W/L_B)^2$ es en general
  mucho menor que $1-\gamma\propto W/L_E$: \textbf{en los transistores modernos $\beta$ está
  limitado por la eficiencia de emisión}, no por la recombinación en la base.
  \item Un PNP tiene peor $\beta$ que un NPN equivalente porque $D_B=D_p<D_n$: los huecos
  difunden más lento. Por eso en circuitos integrados se prefieren los NPN.
\end{itemize}
\end{intuicion}

\begin{ejemplo}[PNP de silicio]
$N_E=\SI{e18}{cm^{-3}}$, $N_B=\SI{e16}{cm^{-3}}$, $N_C=\SI{e15}{cm^{-3}}$, $W=\SI{1}{\micro m}$,
$\tau_E=\SI{0.1}{\micro s}$, $\tau_B=\SI{1}{\micro s}$, $A=\SI{e-4}{cm^2}$, $T=\SI{300}{K}$.
Con las movilidades de la tabla de constantes y Einstein ($D=\mu V_T$):
\[
  D_E = 1350\cdot0.0259 = \SI{35.0}{cm^2/s},\quad
  D_B = 480\cdot0.0259 = \SI{12.4}{cm^2/s},
\]
\[
  L_E=\sqrt{D_E\tau_E}=\SI{18.7}{\micro m},\qquad L_B=\sqrt{D_B\tau_B}=\SI{35.3}{\micro m}.
\]
\[
  1-\gamma\approx\frac{D_E}{D_B}\frac{N_B}{N_E}\frac{W}{L_E}
  =\frac{35.0}{12.4}\cdot\frac{1}{100}\cdot\frac{1}{18.7}=\num{1.51e-3},
  \qquad
  1-\aT\approx\frac12\left(\frac{1}{35.3}\right)^2=\num{4.0e-4}.
\]
\[
  \gamma=0.99849,\quad \aT=0.99960,\quad \adc=0.99809,\quad
  \bdc\approx\frac{1}{\num{1.51e-3}+\num{4.0e-4}}\approx 520 .
\]
Las expresiones exactas dan prácticamente lo mismo ($\bdc=524$). Con
$V_{EB}=\SI{0.65}{V}$ resulta $I_C\approx\SI{1.58}{mA}$ e $I_B\approx\SI{3.0}{\micro A}$.

En un transistor real $\beta$ sería menor: en un emisor tan dopado la movilidad es bastante
menor que la de la tabla (que es para Si poco dopado), y además el gap se achica
(\emph{bandgap narrowing}), lo que aumenta $n_{E0}$ y empeora $\gamma$.
\end{ejemplo}

% =====================================================================
\section{Ecuaciones de Ebers--Moll}
\label{sec:ebers-moll}
% =====================================================================

Sumando las componentes de la sección~\ref{sec:formal}, las corrientes de terminal quedan
\diapo{1.20}:
\[
  I_E = qA\left(\frac{D_E}{L_E}n_{E0}+\frac{D_B}{L_B}\frac{p_{B0}}{\tanh(W/L_B)}\right)\eVEB
      - qA\frac{D_B}{L_B}\frac{p_{B0}}{\sinh(W/L_B)}\eVCB
\]
\[
  I_C = qA\frac{D_B}{L_B}\frac{p_{B0}}{\sinh(W/L_B)}\eVEB
      - qA\left(\frac{D_C}{L_C}n_{C0}+\frac{D_B}{L_B}\frac{p_{B0}}{\tanh(W/L_B)}\right)\eVCB
\]
Tienen siempre la misma estructura: una combinación lineal de $\eVEB$ y $\eVCB$. Ponerles nombre
a los coeficientes da las ecuaciones de Ebers--Moll, válidas en \textbf{las cuatro regiones}.

\begin{clave}[Ecuaciones de Ebers--Moll \enhoja]
\[
  I_E = I_{F0}\eVEB - \aR I_{R0}\eVCB
\]
\[
  I_C = \aF I_{F0}\eVEB - I_{R0}\eVCB
\]
\[
  I_B = I_E - I_C
\]
con
\[
  I_{F0} = qA\left(\frac{D_E}{L_E}n_{E0}+\frac{D_B}{L_B}\frac{p_{B0}}{\tanh(W/L_B)}\right),
  \qquad
  I_{R0} = qA\left(\frac{D_C}{L_C}n_{C0}+\frac{D_B}{L_B}\frac{p_{B0}}{\tanh(W/L_B)}\right),
\]
\[
  \aF I_{F0} = \aR I_{R0} = qA\frac{D_B}{L_B}\frac{p_{B0}}{\sinh(W/L_B)}
  \qquad\text{(relación de reciprocidad)}.
\]
\end{clave}

\paragraph{Significado de cada parámetro.}
\begin{itemize}
  \item $I_{F0}$ ($=I_{ES}$ en otros libros): corriente de saturación de la juntura E--B
  medida con el colector en cortocircuito con la base ($V_{CB}=0$). Con $V_{CB}=0$, $I_E$ es la
  ecuación de un diodo ideal \diapo{1.20}.
  \item $I_{R0}$ ($=I_{CS}$): ídem para la juntura C--B con $V_{EB}=0$.
  \item $\aF$: ganancia de corriente en base común en modo directo ($\aF=I_C/I_E$ con
  $V_{CB}=0$). En MAD coincide con $\adc=\gamma\aT$ \diapo{1.22}.
  \item $\aR$: ídem en modo inverso (intercambiando emisor y colector). Es mucho menor
  que $\aF$ porque el colector está poco dopado (pésima ``eficiencia de emisión inversa'').
\end{itemize}

\begin{figure}[H]
\centering
\begin{circuitikz}[scale=0.95,transform shape]
  \draw (0,0) node[left] {E} to[short,o-] (1,0)
        to[D*,l_=$I_F$,-*] (3,0)
        to[D*,invert,l_=$I_R$,*-] (5,0) to[short,-o] (6,0) node[right] {C};
  \draw (1,0) -- (1,1.6) to[cI,l=$\aR I_R$,invert] (3,1.6) -- (3,0);
  \draw (3,0) -- (3,1.6) to[cI,l=$\aF I_F$] (5,1.6) -- (5,0);
  \draw (3,0) to[short,-o] (3,-1.3) node[below] {B};
  \draw[flecha] (0.2,-0.35) -- node[below,font=\small] {$I_E$} (0.8,-0.35);
  \draw[flecha] (5.2,-0.35) -- node[below,font=\small] {$I_C$} (5.8,-0.35);
  \draw[flecha] (2.7,-0.6) -- node[left,font=\small] {$I_B$} (2.7,-1.1);
\end{circuitikz}
\caption{Circuito equivalente de Ebers--Moll de un PNP (Pierret, fig.~11.6; \diapo{1.19, 1.22}).
$I_F=I_{F0}\eVEB$ e $I_R=I_{R0}\eVCB$ son las corrientes de dos diodos; cada diodo ``alimenta''
una fuente de corriente controlada sobre el otro terminal. El acoplamiento entre los dos diodos
es lo que los diferencia de ``dos diodos enfrentados''.}
\label{fig:em-circuito}
\end{figure}

\begin{figure}[H]
\centering
\begin{circuitikz}[scale=0.95,transform shape]
  \draw (0,0) node[left] {E} to[short,o-] (1,0)
        to[D*,l_=$I_F$,-*] (3,0)
        to[short] (5,0) to[short,-o] (6,0) node[right] {C};
  \draw (3,0) -- (3,1.6) to[cI,l=$\aF I_F$] (5,1.6) -- (5,0);
  \draw (3,0) to[short,-o] (3,-1.3) node[below] {B};
  \node[font=\small,align=left,anchor=west] at (6.8,0.4) {MAD: $I_R\approx-I_{R0}\approx0$\\$I_E\approx I_{F0}e^{V_{EB}/V_T}$\\$I_C\approx\aF I_E$};
\end{circuitikz}
\caption{En activa directa el diodo C--B (en inversa) y la fuente $\aR I_R$ son despreciables:
el transistor es \textbf{un diodo (E--B) más una fuente de corriente controlada} \diapo{1.23}.}
\label{fig:em-mad}
\end{figure}

\begin{ejemplo}[parámetros de Ebers--Moll]
Para el transistor de la sección~\ref{sec:parametros} ($L_C=\sqrt{D_C\tau_C}=\SI{59}{\micro m}$ con
$\tau_C=\SI{1}{\micro s}$):
\[
  I_{F0}=\SI{1.99e-14}{A},\quad I_{R0}=\SI{2.94e-14}{A},\quad
  \aF=0.9981,\quad \aR=\frac{\aF I_{F0}}{I_{R0}}=0.68 .
\]
En modo inverso $\beta_R=\aR/(1-\aR)\approx2$, frente a $\beta_F\approx525$ en modo directo:
el transistor es muy asimétrico.
\end{ejemplo}

\subsection{PNP y NPN}

Todo lo anterior vale para un NPN cambiando huecos por electrones e invirtiendo los signos
de todas las tensiones y corrientes.

\begin{table}[H]
\centering\small
\caption{Traducción PNP $\leftrightarrow$ NPN. Las diapositivas de la 2.\textsuperscript{a} parte
(tomadas de Sedra--Smith) usan NPN.}
\label{tab:pnp-npn}
\begin{tabular}{lll}
\toprule
 & PNP (Pierret, hoja) & NPN (Sedra) \\
\midrule
Portadores útiles en la base & huecos ($p_{B0}=n_i^2/N_B$) & electrones ($n_{B0}=n_i^2/N_A$) \\
Tensión de control en MAD & $V_{EB}>0$ & $V_{BE}>0$ \\
Juntura C--B en MAD & $V_{CB}<0$ & $V_{BC}<0$ \\
Tensión de salida & $V_{EC}>0$ & $V_{CE}>0$ \\
Corrientes ($>0$) & $I_E$ entra; $I_C$, $I_B$ salen & $I_C$, $I_B$ entran; $I_E$ sale \\
Ley exponencial & $I_C=I_Se^{V_{EB}/V_T}$ & $i_C=I_Se^{v_{BE}/V_T}$ \\
\bottomrule
\end{tabular}
\end{table}

% =====================================================================
\section{Características estáticas}
\label{sec:caracteristicas}
% =====================================================================

\subsection{Curva de entrada (emisor común)}

La entrada en emisor común es la juntura E--B: la característica $I_B$ vs.\ $V_{EB}$ (con
$V_{EC}$ fija) es la de un diodo \diapo{1.28}. Como $I_B=I_C/\beta$, la curva $I_C$ vs.\
$V_{EB}$ (característica de \emph{transferencia}) tiene la misma forma, escalada por $\beta$.
En Si conduce apreciablemente a partir de \SIrange{0.5}{0.6}{V}, y a corriente constante
$V_{EB}$ \textbf{disminuye unos \SI{2}{mV/\celsius}} con la temperatura (la curva se corre a la
izquierda al calentarse, porque $n_i^2$ crece).

\begin{ojo}[\diapo{1.28}]
Las curvas mostradas están rotuladas $i_C$ vs.\ $v_{BE}$: son curvas de transferencia, no de
entrada. La forma es la misma (exponencial), así que la conclusión ``es la curva de un
diodo'' vale; pero la corriente de entrada es $i_B$, $\beta$ veces menor.
\end{ojo}

\subsection{Curvas de salida (emisor común)}

Las curvas de salida grafican $I_C$ en función de $V_{EC}$, una curva por cada valor de $I_B$
\diapo{1.29}. Las de la figura~\ref{fig:salida} están calculadas con Ebers--Moll y los
parámetros del ejemplo, agregando el efecto Early.

\begin{figure}[H]
\centering
\begin{tikzpicture}
\begin{axis}[
  width=0.82\textwidth,height=7cm,
  xmin=0,xmax=10,ymin=0,ymax=2.6,
  xlabel={$V_{EC}$ (V)},ylabel={$I_C$ (mA)},
  axis lines=left,grid=major,grid style={gray!20},
  every axis plot/.append style={line width=1.2pt,azul,no markers},
  clip=false,
]
\fill[naranjaclaro,opacity=0.7] (axis cs:0,0) rectangle (axis cs:0.25,2.6);
\fill[gris!15] (axis cs:0,0) rectangle (axis cs:10,0.04);
\addplot table[x=VEC,y=IC1] {datos/salida.dat};
\addplot table[x=VEC,y=IC2] {datos/salida.dat};
\addplot table[x=VEC,y=IC3] {datos/salida.dat};
\addplot table[x=VEC,y=IC4] {datos/salida.dat};
\node[font=\scriptsize,anchor=west] at (axis cs:10,0.60) {$I_B=\SI{1}{\micro A}$};
\node[font=\scriptsize,anchor=west] at (axis cs:10,1.20) {\SI{2}{\micro A}};
\node[font=\scriptsize,anchor=west] at (axis cs:10,1.80) {\SI{3}{\micro A}};
\node[font=\scriptsize,anchor=west] at (axis cs:10,2.40) {\SI{4}{\micro A}};
\node[naranja!70!black,font=\small,anchor=west] at (axis cs:0.3,0.25) {$\leftarrow$ saturación};
\node[font=\small] at (axis cs:5,2.05) {activa directa};
\node[gris,font=\scriptsize,anchor=south west] at (axis cs:3,0.04) {corte ($I_B=0$)};
\end{axis}
\end{tikzpicture}
\caption{Curvas de salida de un PNP en emisor común (calculadas). En activa las curvas están
equiespaciadas ($\Delta I_C=\beta\Delta I_B$) y levemente inclinadas (efecto Early). Por
debajo de $V_{EC}\approx\SIrange{0.1}{0.2}{V}$ la juntura C--B entra en directa y el
transistor se satura: $I_C$ cae rápidamente a cero.}
\label{fig:salida}
\end{figure}

\paragraph{Base común.} Si se grafica $I_C$ vs.\ $V_{CB}$ con $I_E$ como parámetro, las curvas
son prácticamente horizontales ($I_C\approx\adc I_E$), sin efecto Early visible, y siguen
así hasta \emph{pasar} $V_{CB}=0$: la saturación empieza recién cuando la juntura C--B entra
claramente en directa ($V_{CB}\approx\SI{0.5}{V}$).

\subsection{Corrientes de fuga: $I_{CB0}$ e $I_{CE0}$}

Incluyendo la corriente inversa de la juntura C--B \diapo{2.9--2.11}:

\begin{clave}[Base común y emisor común \enhoja]
\[
  I_C = \adc I_E + I_{CB0}
  \qquad\Longrightarrow\qquad
  I_C = \frac{\adc}{1-\adc}I_B + \frac{I_{CB0}}{1-\adc} = \bdc I_B + I_{CE0},
\]
\[
  I_{CE0} = \frac{I_{CB0}}{1-\adc} = (\bdc+1)\,I_{CB0}.
\]
\begin{itemize}[nosep]
  \item $I_{CB0}$: corriente de colector con el \textbf{emisor abierto} ($I_E=0$).
  \item $I_{CE0}$: corriente de colector con la \textbf{base abierta} ($I_B=0$).
\end{itemize}
\end{clave}

\begin{intuicion}[por qué $I_{CE0}\gg I_{CB0}$]
Con la base abierta, los electrones de la corriente de fuga C--B que llegan a la base no
tienen por dónde salir: se acumulan, polarizan en directa a la juntura E--B y el emisor
inyecta huecos. Esa fuga actúa como una corriente de base y el transistor la
\emph{amplifica} por $\beta+1$. Por eso, y porque $I_{CB0}$ se duplica cada \SI{10}{\celsius}
aproximadamente, la fuga es un problema en circuitos analógicos que hay que compensar
\diapo{2.34}.
\end{intuicion}

\begin{ojo}[$I_{CB0}$ \diapo{2.11}]
``En nuestro desarrollo sería la $I_{Cn}$'': es una aproximación. Con Ebers--Moll, emisor
abierto ($I_E=0$) y $V_{CB}\ll-V_T$:
\[
  I_{CB0} = I_{R0}\left(1-\aF\aR\right),
\]
que incluye $I_{Cn}$ más la pequeña corriente de huecos de la base hacia el colector. Además
hay que sumarle la corriente de generación en la ZCE de C--B, que en Si suele dominar.
\end{ojo}

\subsection{Gummel plot y dependencia de $\beta$}

$\beta$ es la figura de mérito más usada para comparar transistores, pero \textbf{no es
constante}: depende de $I_C$ y de la temperatura \diapo{1.27}. El \emph{gráfico de
Gummel} muestra $\log I_C$ y $\log I_B$ en función de $V_{EB}$; la distancia vertical entre
ambas curvas es $\log\beta$.

\begin{figure}[H]
\centering
\begin{tikzpicture}
\begin{semilogyaxis}[
  name=gum,
  width=0.47\textwidth,height=6.2cm,
  xmin=0.2,xmax=0.95,ymin=1e-13,ymax=1e-1,
  xlabel={$V_{EB}$ (V)},ylabel={corriente (A)},
  legend pos=north west,legend style={font=\scriptsize},
  every axis plot/.append style={line width=1.1pt},
  ytick={1e-12,1e-10,1e-8,1e-6,1e-4,1e-2},
]
\addplot[azul] table[x=VEB,y=IC] {datos/gummel.dat};
\addlegendentry{$I_C$}
\addplot[rojo] table[x=VEB,y=IB] {datos/gummel.dat};
\addlegendentry{$I_B$}
\draw[<->,gris] (axis cs:0.6,2.2e-4) -- node[right,font=\scriptsize,black] {$\beta$} (axis cs:0.6,7e-7);
\node[font=\scriptsize,align=center,rojo] at (axis cs:0.36,2e-8) {R--G en la\\ZCE E--B};
\node[font=\scriptsize,align=center,azul] at (axis cs:0.83,1e-6) {alta\\inyección};
\end{semilogyaxis}
\begin{axis}[
  at={(gum.east)},anchor=west,xshift=1.4cm,
  width=0.47\textwidth,height=6.2cm,
  xmode=log,xmin=1e-6,xmax=1e-1,ymin=0,ymax=360,
  xlabel={$I_C$ (A)},ylabel={$\beta$},
  every axis plot/.append style={line width=1.1pt,smooth},
  legend pos=north west,legend style={font=\scriptsize},
]
\addplot[rojo] coordinates {(1e-6,150)(1e-5,215)(1e-4,270)(1e-3,305)(1e-2,275)(3e-2,200)(1e-1,95)};
\addlegendentry{\SI{125}{\celsius}}
\addplot[azul] coordinates {(1e-6,100)(1e-5,145)(1e-4,180)(1e-3,195)(1e-2,170)(3e-2,120)(1e-1,55)};
\addlegendentry{\SI{25}{\celsius}}
\addplot[violeta] coordinates {(1e-6,65)(1e-5,90)(1e-4,105)(1e-3,110)(1e-2,90)(3e-2,60)(1e-1,25)};
\addlegendentry{\SI{-55}{\celsius}}
\end{axis}
\end{tikzpicture}
\caption{Izquierda: gráfico de Gummel (Pierret, fig.~11.13). Derecha: $\beta$ en función de la
corriente de colector y la temperatura (valores típicos, \diapo{1.27}). El modelo ideal sólo
vale en el tramo central donde las dos rectas son paralelas.}
\label{fig:gummel}
\end{figure}

\begin{itemize}
  \item \textbf{Corrientes bajas:} a $I_B$ se le suma la corriente de recombinación en la ZCE
  de E--B, que crece como $e^{V_{EB}/2V_T}$ (pendiente mitad). $I_B$ ya no es proporcional a
  $I_C$ y $\beta$ cae.
  \item \textbf{Corrientes altas:} alta inyección en la base (la pendiente de $I_C$ pasa a
  $e^{V_{EB}/2V_T}$), aglomeración de corriente en los bordes del emisor
  (\emph{current crowding}) y caídas en las resistencias serie. $\beta$ vuelve a caer.
  \item \textbf{Temperatura:} $\beta$ aumenta con $T$ (sobre todo porque mejora la eficiencia
  de emisión).
\end{itemize}

Por eso, al decir ``el transistor tiene $\beta=200$'' hay que aclarar a qué corriente y
temperatura. En la materia se toma $\beta$ constante dentro del rango de validez del modelo
\diapo{1.31}.

% =====================================================================
\section{Apartamientos del modelo ideal}
% =====================================================================

\subsection{Modulación del ancho de base (efecto Early)}
\label{sec:early}

Al aumentar la tensión inversa $|V_{CB}|$, la ZCE de la juntura C--B se ensancha y una parte
de ese ensanchamiento invade la base: el ancho cuasineutro
\[
  W = W_B - x_{nEB}(V_{EB}) - x_{nCB}(V_{CB})
\]
\textbf{disminuye} \diapo{2.3}. Como en MAD
\[
  I_C \simeq qA\frac{D_B}{W}\,p_{B0}\,e^{V_{EB}/V_T}
  \qquad\left(\text{y también } \beta\simeq\frac{1}{\frac{D_E}{D_B}\frac{N_B}{N_E}\frac{W}{L_E}+\frac12\frac{W^2}{L_B^2}}\right),
\]
$I_C$ y $\beta$ \textbf{crecen} con $|V_{CB}|$ aunque $V_{EB}$ se mantenga fija. En las curvas de
salida eso se ve como una pendiente: las rectas de la región activa, prolongadas hacia la
izquierda, se cortan aproximadamente en un mismo punto del eje, $V_{EC}=-V_A$, donde $V_A$ es
la \textbf{tensión de Early} \diapo{2.4}.

\begin{figure}[H]
\centering
\begin{tikzpicture}
\begin{axis}[
  width=0.8\textwidth,height=6cm,
  xmin=-45,xmax=12,ymin=0,ymax=3.2,
  axis lines=middle,
  xlabel={$V_{EC}$},ylabel={$I_C$},
  xlabel style={at={(ticklabel* cs:1)},anchor=north west},
  ytick=\empty,xtick={-40},xticklabels={$-V_A$},
  every axis plot/.append style={line width=1.2pt},
]
\foreach \k in {0.6,1.2,1.8,2.4}{
  \edef\tmp{\noexpand\addplot[azul,domain=0.4:12] {\k*(1+x/40)};
            \noexpand\addplot[azul,dashed,thin,domain=-40:0.4] {\k*(1+x/40)};}
  \tmp
}
\fill[azul] (axis cs:-40,0) circle (2pt);
\node[font=\scriptsize,anchor=west,align=left] at (axis cs:-36,2.6) {pendiente $\dfrac{\partial I_C}{\partial V_{EC}}\simeq\dfrac{I_C}{V_A}=\dfrac{1}{r_o}$};
\end{axis}
\end{tikzpicture}
\caption{Efecto Early: las curvas de salida extrapoladas convergen en $-V_A$ (Sedra, fig.~6.18;
\diapo{2.4}). Valores típicos: $V_A=\SIrange{20}{200}{V}$.}
\label{fig:early}
\end{figure}

\begin{clave}[Efecto Early]
\[
  I_C \simeq I_S\,e^{V_{EB}/V_T}\left(1+\frac{V_{EC}}{V_A}\right),
  \qquad
  r_o = \left(\frac{\partial I_C}{\partial V_{EC}}\right)^{-1} = \frac{V_A}{I_C}
  \ \text{\enhojatexto},
  \quad
  V_A \approx \frac{W}{\left|\dd W/\dd V_{CB}\right|}.
\]
\end{clave}

La salida del transistor no es una fuente de corriente ideal: tiene una resistencia en paralelo
$r_o$, finita. Para tener $V_A$ grande conviene que la ZCE de C--B penetre poco en la base, es
decir \textbf{$N_C<N_B$} (la ZCE se extiende sobre todo hacia el colector) y una base no
demasiado angosta. Hay un compromiso con $\beta$: base angosta = $\beta$ alto pero $V_A$ bajo.

\begin{ejemplo}[tensión de Early]
Para el transistor del ejemplo ($N_B=\SI{e16}{}$, $N_C=\SI{e15}{cm^{-3}}$, $W=\SI{1}{\micro m}$ con
$V_{CB}=\SI{-5}{V}$). La parte de la ZCE de C--B que cae en la base es
\[
  x_{nCB} = \sqrt{\frac{2\epsilon_r\epsilon_0}{q}\,\frac{N_C}{N_B(N_B+N_C)}\,(V_{bi}-V_{CB})},
  \qquad V_{bi}=V_T\ln\frac{N_BN_C}{n_i^2}=\SI{0.656}{V}.
\]
Con $\epsilon_r=11.8$: $x_{nCB}(\SI{-5}{V})=\SI{0.259}{\micro m}$ y
$\dd x_{nCB}/\dd|V_{CB}|=x_{nCB}/[2(V_{bi}-V_{CB})]=\SI{0.0229}{\micro m/V}$, así que
\[
  V_A\approx\frac{\SI{1}{\micro m}}{\SI{0.0229}{\micro m/V}}\approx\SI{44}{V},
  \qquad
  r_o=\frac{V_A}{I_C}=\frac{\SI{44}{V}}{\SI{1.58}{mA}}\approx\SI{28}{k\ohm}.
\]
\end{ejemplo}

\begin{ojo}[\diapo{2.3}]
``Asumiendo $V_{CB}>5V_T$'': en un PNP en MAD la juntura C--B está en inversa, o sea
$V_{CB}<0$. La condición correcta es $V_{CB}<-5V_T$ (es decir, $|V_{CB}|\gtrsim\SI{0.13}{V}$),
que permite despreciar $e^{V_{CB}/V_T}$ frente a 1.
\end{ojo}

\subsection{Perforación (\emph{punch-through})}

Si $|V_{CB}|$ sigue aumentando, la ZCE de C--B puede llegar a tocar la de E--B: $W\to0$
\diapo{2.5--2.6}. A partir de ahí la tensión de colector actúa directamente sobre la barrera
de la juntura E--B, la baja, y el emisor inyecta una corriente enorme que ya no controla la
base. La corriente queda limitada sólo por el circuito externo; si la potencia excede el
área de operación segura (SOA) el daño es irreversible.

La tensión de perforación $V_{PT}$ se obtiene de $W_B - x_{nEB} - x_{nCB}(V_{PT}) = 0$.

\begin{ejemplo}[¿perforación o avalancha?]
En el transistor del ejemplo, $W_B = W + x_{nEB}+x_{nCB} = 1 + 0.155 + 0.259 = \SI{1.41}{\micro m}$
($x_{nEB}=\SI{0.155}{\micro m}$ con $V_{EB}=\SI{0.65}{V}$, $V_{bi,EB}=\SI{0.835}{V}$). La perforación
ocurre cuando $x_{nCB}=W_B-x_{nEB}=\SI{1.26}{\micro m}$:
\[
  V_{bi}-V_{PT} = \frac{(\SI{1.26e-4}{cm})^2}{\frac{2\epsilon_r\epsilon_0}{q}\frac{N_C}{N_B(N_B+N_C)}}
  \approx\SI{134}{V}
  \quad\Rightarrow\quad V_{CB,PT}\approx\SI{-133}{V}.
\]
Una juntura de Si con $N_C=\SI{e15}{cm^{-3}}$ tiene una ruptura por avalancha de unos
\SI{300}{V}, así que en este transistor la perforación ocurriría antes. Para evitarla se
dopa más la base o se la hace más ancha (a costa de $\beta$).
\end{ejemplo}

\subsection{Ruptura por avalancha}

La juntura C--B es la que soporta la tensión alta \diapo{2.7}. Si el campo es suficientemente
grande, los portadores que la cruzan generan pares por \textbf{ionización por impacto} y la
corriente se multiplica por
\[
  M = \frac{1}{1-\left(|V_{CB}|/BV_{CB0}\right)^n}\qquad (2<n<10),
\]
donde $BV_{CB0}$ es la tensión de ruptura de la juntura C--B con el emisor abierto.

\paragraph{Base común (emisor abierto o $I_E$ fijo).} $I_C=M(\adc I_E+I_{CB0})$: la corriente
sólo diverge cuando $M\to\infty$, es decir en $|V_{CB}|=BV_{CB0}$.

\paragraph{Emisor común con base abierta ($I_B=0$).} Ahora $I_E=I_C$, y
\[
  I_C = M(\adc I_C+I_{CB0})\quad\Rightarrow\quad I_C=\frac{M\,I_{CB0}}{1-M\adc},
\]
que diverge cuando $M\adc=1$, mucho antes de que $M\to\infty$:
\begin{clave}[Tensión de ruptura en emisor común]
\[
  M\adc=1 \;\Rightarrow\; 1-\left(\frac{BV_{CE0}}{BV_{CB0}}\right)^n=\adc
  \;\Rightarrow\;
  BV_{CE0} = BV_{CB0}\,(1-\adc)^{1/n} \simeq \frac{BV_{CB0}}{\sqrt[n]{\beta}} .
\]
\end{clave}

\begin{intuicion}[por qué la ruptura en emisor común es más baja]
Es el mismo mecanismo que hace $I_{CE0}=(\beta+1)I_{CB0}$: en un PNP, los electrones generados
por avalancha en la ZCE de C--B son arrastrados a la base; con la base abierta no tienen salida,
polarizan más a la juntura E--B y el emisor inyecta más huecos, que generan más pares\ldots{}
Es una realimentación positiva con ganancia $\beta$, que dispara la corriente con una
multiplicación pequeña ($M=1/\adc\approx1.01$).
\end{intuicion}

\begin{figure}[H]
\centering
\begin{tikzpicture}
\begin{axis}[
  width=0.75\textwidth,height=5.8cm,
  xmin=0,xmax=105,ymin=0,ymax=5,restrict y to domain=0:6,
  xlabel={$V_{EC}$ (V)},ylabel={$I_C$},
  axis lines=left,ytick=\empty,
  xtick={52,100},xticklabels={$BV_{CE0}$,$BV_{CB0}$},
  every axis plot/.append style={line width=1.1pt,azul,samples=120},
]
\foreach \k in {0.0,0.8,1.6,2.4}{
  \edef\tmp{\noexpand\addplot[domain=0.5:51.6] {(\k+0.08)*(1+x/150)/(1-(x/52)^7)};}
  \tmp
}
\draw[gris,dashed] (axis cs:52,0) -- (axis cs:52,5);
\draw[gris,dashed] (axis cs:100,0) -- (axis cs:100,5);
\node[font=\scriptsize,anchor=west] at (axis cs:2,0.35) {$I_B=0$};
\end{axis}
\end{tikzpicture}
\caption{Ruptura en emisor común (\diapo{2.8}, Pierret fig.~11.10). Todas las curvas se
disparan cerca de $BV_{CE0}$, bastante antes que la ruptura de la juntura aislada $BV_{CB0}$.
Con $\beta=100$ y $n\approx7$ resulta $BV_{CE0}\approx0.52\,BV_{CB0}$.}
\label{fig:avalancha}
\end{figure}

\begin{ojo}[``$BV_{CE0}=0.52\,BV_{CB0}$'' \diapo{2.8}]
No es una regla general: corresponde a un caso particular ($\beta\approx100$ y $n\approx7$,
porque $100^{-1/7}=0.52$). Con $n=4$ y el mismo $\beta$ daría $0.32$. La relación general es
$BV_{CE0}\simeq BV_{CB0}\,\beta^{-1/n}$.
\end{ojo}

\subsection{Fototransistor}

Si la luz incide sobre la ZCE de C--B, los pares fotogenerados son separados por el campo: los
huecos van al colector y los electrones a la base \diapo{2.7}. Esa fotocorriente $I_L$ cumple
el papel de una corriente de base (aunque el terminal B esté abierto) y el transistor la
amplifica:
\[
  I_C = (\beta+1)\,I_L .
\]
El fototransistor es mucho más sensible que un fotodiodo, pero también mucho más lento (la
ganancia se paga con ancho de banda, como en el fotoconductor de la unidad 6).

\subsection{Otras desviaciones}

\begin{itemize}
  \item \textbf{Recombinación en la ZCE de E--B}: degrada $\beta$ a corrientes bajas (Gummel).
  \item \textbf{Alta inyección y resistencia de base}: degradan $\beta$ a corrientes altas.
  \item \textbf{Base gradual}: en los transistores reales la base se dopa por difusión o
  implantación y $N_B$ no es uniforme: aparece un campo interno que \emph{acelera} a los
  minoritarios, mejora $\aT$ y la velocidad.
  \item \textbf{Geometría real}: el ``efecto transistor'' ocurre sólo en el pequeño volumen
  vertical debajo del emisor \diapo{2.23}; el resto agrega resistencias ($r_b$, $r_c$) y
  capacidades parásitas.
\end{itemize}

% =====================================================================
\section{Saturación y corte: el transistor como llave}
% =====================================================================

Para usar al transistor como llave hacen falta dos estados \diapo{2.12}: \textbf{abierta}
($I\approx0$) y \textbf{cerrada} ($V\approx0$). Corresponden a las regiones de corte y de
saturación.

\begin{itemize}
  \item \textbf{Corte} (ambas junturas en inversa): sólo circulan corrientes de fuga del
  orden de $I_{CB0}$ e $I_{CE0}$. La llave está abierta.
  \item \textbf{Saturación} (ambas junturas en directa): $V_{EC}$ es muy chica
  (\SIrange{0.05}{0.2}{V}) y la corriente la fija el circuito externo. La llave está cerrada.
\end{itemize}

\subsection{¿Cómo puede $V_{EC}$ ser menor que $V_{EB}$?}

\diapo{2.12} pregunta cómo puede ser $V_{EC}<\SI{0.7}{V}$ si $V_{EB}\approx\SI{0.7}{V}$. La
respuesta está en $V_{EC}=V_{EB}-V_{CB}$: en saturación \emph{ambas} junturas están en directa y
sus tensiones casi se cancelan. Por ejemplo, $V_{EB}=\SI{0.64}{V}$ y $V_{CB}=\SI{0.53}{V}$ dan
$V_{EC}=\SI{0.11}{V}$.

Con Ebers--Moll, imponiendo $I_C$ e $I_B$ (las ecuaciones son lineales en $e^{V_{EB}/V_T}$ y $e^{V_{CB}/V_T}$, así que se despejan fácilmente) y despreciando los ``$-1$'', se obtiene:

\begin{clave}[Tensión de saturación]
\[
  V_{EC,sat} = V_T\ln\left[\frac{\dfrac{1}{\aR}+\dfrac{\beta_{forz}}{\beta_R}}
                               {1-\dfrac{\beta_{forz}}{\beta_F}}\right],
  \qquad
  \beta_{forz}=\frac{I_C}{I_B}<\beta_F,
  \qquad \beta_R=\frac{\aR}{1-\aR}.
\]
\end{clave}

\begin{ejemplo}[$V_{EC,sat}$]
Con el transistor del ejemplo ($\beta_F=525$, $\aR=0.68$, $\beta_R=2.1$) y $I_C=\SI{1}{mA}$:
\begin{center}
\begin{tabular}{cccc}
\toprule
$I_B$ & $\beta_{forz}$ & $V_{EC,sat}$ & comentario \\
\midrule
\SI{5}{\micro A}   & 200 & \SI{131}{mV} & saturación leve \\
\SI{10}{\micro A}  & 100 & \SI{106}{mV} & \\
\SI{100}{\micro A} & 10  & \SI{48}{mV}  & saturación profunda \\
\bottomrule
\end{tabular}
\end{center}
Cuanto más se ``sobreexcita'' la base, más se acerca a cero la tensión de la llave.
(Verificado resolviendo numéricamente Ebers--Moll.)
\end{ejemplo}

\subsection{Recta de carga y transferencia}

\begin{figure}[H]
\centering
\begin{minipage}{0.38\textwidth}
\centering
\begin{circuitikz}[scale=0.85,transform shape]
  \draw (0,0) node[pnp](Q){};
  \draw (Q.E) -- ++(0,0.5) node[vcc]{$V_{CC}=\SI{5}{V}$};
  \draw (Q.C) to[R,l=$R_C$,-*] ++(0,-2) coordinate(o) -- ++(0,-0.6) node[ground]{};
  \draw (o) to[short,-o] ++(1.2,0) node[right] {$v_o$};
  \draw (Q.B) to[short,-o] ++(-0.8,0) node[left] {$v_{in}$};
  \node[font=\scriptsize,align=center] at (-1.3,-1.2) {$V_{EB}=V_{CC}-v_{in}$};
\end{circuitikz}
\end{minipage}%
\begin{minipage}{0.6\textwidth}
\centering
\begin{tikzpicture}
\begin{axis}[
  width=\textwidth,height=5.5cm,
  xmin=0.4,xmax=0.75,ymin=0,ymax=5.5,
  xlabel={$V_{EB}$ (V)},ylabel={$v_o=I_CR_C$ (V)},
  axis lines=left,grid=major,grid style={gray!20},
]
\fill[gris!15] (axis cs:0.4,0) rectangle (axis cs:0.58,5.5);
\fill[naranjaclaro] (axis cs:0.682,0) rectangle (axis cs:0.75,5.5);
\addplot[azul,line width=1.3pt] table[x=VEB,y=vo] {datos/transfer.dat};
\node[font=\small] at (axis cs:0.49,4.5) {corte};
\node[font=\small] at (axis cs:0.635,4.5) {activa};
\node[font=\small,rotate=90] at (axis cs:0.715,2.5) {saturación};
\end{axis}
\end{tikzpicture}
\end{minipage}
\caption{Inversor con PNP ($R_C=\SI{1}{k\ohm}$) y su característica de transferencia, calculada
con Ebers--Moll (equivalente PNP de la \diapo{2.14}). En corte $v_o\approx0$; en saturación
$v_o\approx V_{CC}-V_{EC,sat}$. La región activa es la zona de pendiente alta: allí el circuito
amplifica.}
\label{fig:llave}
\end{figure}

\paragraph{La malla de salida decide la saturación \diapo{2.14}.} Con $V_{EB}$ fija en activa,
$I_C=\beta I_B$ es independiente de $R_C$. Pero la malla de colector impone
$V_{EC}=V_{CC}-I_CR_C$: si $\beta I_BR_C$ supera $V_{CC}$, la tensión $V_{EC}$ ``no alcanza'',
la juntura C--B entra en directa y el transistor se satura. La corriente queda limitada a
$I_{C,sat}\approx(V_{CC}-V_{EC,sat})/R_C<\beta I_B$. Criterio práctico de diseño: para
asegurar saturación, $I_B>I_{C,sat}/\beta_{\min}$ con un margen.

\subsection{Qué pasa con $\beta$ en saturación}

En saturación las dos junturas inyectan: además del emisor, el colector inyecta huecos hacia
la base y la base inyecta electrones hacia el colector. La corriente de colector neta
\emph{baja} (las dos inyecciones se restan) y la corriente de base \emph{sube} (hay que
alimentar las dos recombinaciones). Resultado: $\beta_{forz}=I_C/I_B$ cae mucho
\diapo{2.15}. En circuitos lineales se evita a toda costa; en conmutación se la busca.

\begin{intuicion}[carga almacenada y tiempo de almacenamiento]
En saturación hay un exceso grande de huecos en toda la base (el perfil de la
figura~\ref{fig:perfiles} ya no llega a cero en $x=W$) y también en el colector. Para apagar la
llave primero hay que \emph{retirar esa carga}: mientras tanto el transistor sigue conduciendo.
Ese retardo (\emph{storage time}) limita la velocidad de conmutación; por eso los circuitos
rápidos evitan la saturación profunda (p.\,ej.\ con un diodo Schottky entre base y colector).
\end{intuicion}

% =====================================================================
\section{Modelos de pequeña señal}
\label{sec:pequena-senal}
% =====================================================================

\subsection{Linealización}

En un amplificador el transistor se polariza en un punto de reposo $Q$ (en MAD) y se le
superpone una señal chica \diapo{2.16--2.17}. Notación: mayúscula con subíndice mayúscula
para continua ($I_C$), minúscula con subíndice minúscula para la señal ($i_c$) y minúscula con
subíndice mayúscula para el total ($i_C=I_C+i_c$).

Si la señal es suficientemente chica se puede desarrollar en serie de Taylor a primer orden
alrededor de $Q$ (emisor común: entrada $v_{be}$, salida $v_{ce}$):
\[
  i_b = \left.\frac{\partial I_B}{\partial V_{BE}}\right|_{Q}v_{be}
      + \left.\frac{\partial I_B}{\partial V_{CE}}\right|_{Q}v_{ce}
      = y_{11}v_{be}+y_{12}v_{ce},
  \qquad
  i_c = y_{21}v_{be}+y_{22}v_{ce}.
\]
Todas las derivadas tienen unidades de A/V: son \textbf{admitancias} \diapo{2.18--2.19}.

\begin{intuicion}[¿qué tan chica es ``pequeña señal''?]
Como $i_C=I_Ce^{v_{be}/V_T}=I_C(1+v_{be}/V_T+\tfrac12(v_{be}/V_T)^2+\dots)$, el término lineal
domina si $v_{be}\ll V_T$. Con $v_{be}=\SI{5}{mV}$ el término cuadrático ya es el
\SI{10}{\percent} del lineal: en la práctica ``pequeña señal'' significa unos pocos mV en la
juntura base--emisor.
\end{intuicion}

\subsection{Modelo híbrido-\texorpdfstring{$\pi$}{pi} de Giacoletto}

Se reescriben los parámetros $y$ con significado físico \diapo{2.20}:

\begin{clave}[Parámetros del modelo híbrido-$\pi$ \enhoja]
\[
  g_m = \left.\frac{\partial I_C}{\partial V_{EB}}\right|_Q = \frac{I_C}{V_T}
  \qquad(y_{21},\ \text{transconductancia})
\]
\[
  r_\pi = \left(\left.\frac{\partial I_B}{\partial V_{EB}}\right|_Q\right)^{-1}
        = \frac{V_T}{I_B} = \frac{\beta}{g_m}
  \qquad(1/y_{11})
\]
\[
  r_o = \frac{V_A}{I_C}\quad(1/y_{22},\ \text{efecto Early}),
  \qquad
  r_\mu\approx\beta\,r_o,
  \qquad
  g_mr_o = \frac{V_A}{V_T}.
\]
$y_{12}$ (realimentación de la salida sobre la entrada) es despreciable en primera aproximación;
el efecto residual se representa con $r_\mu$.
\end{clave}

\begin{figure}[H]
\centering
\begin{circuitikz}[scale=0.95,transform shape]
  % entrada
  \draw (0,3) node[left] {B} to[short,o-] (0.3,3) to[R] (2,3);
  \node[above,font=\small] at (1.15,3.35) {$r_b$};
  \draw (2,3) to[short,*-] (2,3) node[above left,font=\small] {B'};
  \draw (2,3) to[C] (2,0);
  \node[left,font=\small] at (1.6,1.5) {$C_\pi$};
  \draw (2,3) -- (3.6,3) to[R] (3.6,0);
  \node[right,font=\small] at (3.95,1.5) {$r_\pi$};
  \node[font=\small,azul] at (3.0,2.5) {$+$}; \node[font=\small,azul] at (3.0,0.5) {$-$};
  \node[font=\small,azul] at (3.0,1.5) {$v_\pi$};
  % realimentacion
  \draw (2,3) -- (2,4.2) -- (3.4,4.2) to[R] (6.4,4.2) -- (7.2,4.2);
  \node[above,font=\small] at (4.9,4.55) {$r_\mu$};
  \draw (2,4.2) -- (2,5.4) -- (3.4,5.4) to[C] (6.4,5.4) -- (7.2,5.4) -- (7.2,3);
  \node[above,font=\small] at (4.9,5.85) {$C_\mu$};
  \draw (7.2,4.2) to[short,*-] (7.2,4.2);
  % salida
  \draw (3.6,3) to[short,*-] (3.6,3);
  \draw (7.2,3) to[short,*-] (7.2,3) to[R] (7.2,0);
  \node[left,font=\small] at (6.85,1.5) {$r_o$};
  \draw (7.2,3) -- (8.8,3) to[cI] (8.8,0);
  \node[right,font=\small] at (9.3,1.5) {$g_mv_\pi$};
  \draw (8.8,3) to[short,-o] (10.6,3) node[right] {C};
  \draw (0.3,0) node[left] {E} to[short,o-] (10.6,0) to[short,-o] (10.6,0) node[right] {E};
  \foreach \x in {2,3.6,7.2,8.8} \fill (\x,0) circle (1.6pt);
\end{circuitikz}
\caption{Modelo híbrido-$\pi$ completo (Giacoletto) \diapo{2.22}. $r_b$ es la resistencia
distribuida entre el contacto de base y la base ``activa'' B' (debajo del emisor). El mismo
modelo vale para PNP y NPN, porque relaciona \emph{variaciones}.}
\label{fig:giacoletto}
\end{figure}

\begin{figure}[H]
\centering
\begin{circuitikz}[scale=0.9,transform shape]
  \draw (0,2.4) node[left] {B} to[short,o-] (1.5,2.4) to[R] (1.5,0);
  \node[right,font=\small] at (1.85,1.2) {$r_\pi$};
  \node[font=\small,azul] at (0.7,1.9) {$+$}; \node[font=\small,azul] at (0.7,0.5) {$-$};
  \node[font=\small,azul] at (0.7,1.2) {$v_{be}$};
  \draw (1.5,2.4) -- (1.5,2.4);
  \draw (4.6,2.4) to[cI] (4.6,0);
  \node[left,font=\small] at (4.1,1.2) {$g_mv_{be}$};
  \draw (4.6,2.4) -- (6.4,2.4) to[R] (6.4,0);
  \node[right,font=\small] at (6.75,1.2) {$r_o$};
  \draw (6.4,2.4) to[short,-o] (7.6,2.4) node[right] {C};
  \draw (0,0) node[left] {E} to[short,o-] (7.6,0) to[short,-o] (7.6,0);
  \node[font=\small,align=left,anchor=west] at (8.6,1.4) {$i_b=v_{be}/r_\pi$\\[2pt]$i_c=g_mv_{be}+v_{ce}/r_o$};
\end{circuitikz}
\caption{Modelo simplificado de baja frecuencia \diapo{2.20}: el que se usa para casi todos
los cálculos de ganancia.}
\label{fig:pi-simple}
\end{figure}

\begin{ejemplo}[parámetros en un punto de reposo]
$I_C=\SI{1}{mA}$, $\beta=200$, $V_A=\SI{50}{V}$, $V_T=\SI{26}{mV}$:
\[
  g_m=\frac{\SI{1}{mA}}{\SI{26}{mV}}=\SI{38.5}{mS},\quad
  r_\pi=\frac{200}{\SI{38.5}{mS}}=\SI{5.2}{k\ohm},\quad
  r_o=\frac{\SI{50}{V}}{\SI{1}{mA}}=\SI{50}{k\ohm},\quad
  r_\mu\approx\SI{10}{M\ohm}.
\]
Ganancia de tensión intrínseca máxima (emisor común sin carga):
$|A_v|=g_mr_o=V_A/V_T\approx1900$. Notar que $g_m$ depende \emph{sólo} de la corriente de
polarización, no del transistor: es una propiedad de la ley exponencial.
\end{ejemplo}

\subsection{Capacidades internas}

Cada juntura tiene una capacidad de juntura (carga en la ZCE) y una de difusión (carga de
minoritarios en exceso). En MAD \diapo{2.21}:
\begin{itemize}
  \item \textbf{Juntura C--B (inversa):} sólo capacidad de juntura,
  $C_\mu(V_{CB}) = \dfrac{C_{\mu0}}{\left(1-V_{CB}/V_{bi}\right)^{n_c}}$, con $n_c\approx1/2$
  (abrupta) a $1/3$ (gradual).
  \item \textbf{Juntura E--B (directa):} capacidad de juntura $C_{je}$ (en directa la fórmula
  diverge; se usa la aproximación $C_{je}\approx2C_{je0}$) más la capacidad de difusión de la
  base, $C_b$.
\end{itemize}

\paragraph{Capacidad de difusión de la base.} La carga de huecos en exceso en la base (perfil
triangular, figura~\ref{fig:perfiles}) es
\[
  Q_B = qA\,\frac{\Delta p_B(0)\,W}{2}
  \qquad\text{y}\qquad
  I_C = qA\,\frac{D_B\,\Delta p_B(0)}{W}
  \quad\Longrightarrow\quad
  Q_B = \tau_F\,I_C,\qquad \tau_F = \frac{W^2}{2D_B}.
\]
$\tau_F$ es el \textbf{tiempo de tránsito} por la base: el tiempo que tarda en promedio un
hueco en cruzarla. Derivando respecto de $V_{EB}$:
\begin{clave}[Capacidades del modelo \enhoja]
\[
  C_b = \frac{\dd Q_B}{\dd V_{EB}} = \tau_F\,\frac{\dd I_C}{\dd V_{EB}} = \tau_F\,g_m,
  \qquad
  C_\pi = C_b + C_{je} \approx \tau_F\frac{I_C}{V_T} + 2C_{je0}.
\]
\end{clave}

\begin{ojo}[notación de la hoja y \diapo{2.21}]
\begin{itemize}
  \item En la hoja de fórmulas $C_b=\tau_B\,g_m$: ahí \textbf{$\tau_B$ es el tiempo de tránsito
  por la base} ($=\tau_F=W^2/2D_B$), \emph{no} la vida media de los huecos en la base que
  aparece en $L_B=\sqrt{D_B\tau_B}$. Son cantidades muy distintas: en el ejemplo, la vida media
  es \SI{1}{\micro s} y el tránsito \SI{0.4}{ns}. Su cociente es justamente
  $\tau_B^{vida}/\tau_F=2L_B^2/W^2\approx1/(1-\aT)$.
  \item La expresión $C_{De}=\frac{qA}{V_T}(L_En_{E0}+L_Bp_{B0})e^{V_{EB}/V_T}$ de la diapositiva
  es la de un \emph{diodo largo}; en la base angosta de un transistor el término de la base
  se reemplaza por $\frac{qA}{V_T}\frac{W}{2}p_{B0}e^{V_{EB}/V_T}=\tau_Fg_m$, que es el que
  usa la propia diapo más abajo. Lo mismo para $C_{Dc}$, que en MAD es despreciable.
\end{itemize}
\end{ojo}

\subsection{Respuesta en frecuencia: \texorpdfstring{$f_\beta$}{f-beta} y \texorpdfstring{$f_T$}{fT}}

Con la salida en cortocircuito ($v_{ce}=0$) y despreciando $r_b$ y $r_\mu$, la corriente de base
se reparte entre $r_\pi$ y $C_\pi+C_\mu$:
\[
  v_{be} = \frac{i_b}{1/r_\pi+j\omega(C_\pi+C_\mu)},
  \qquad
  i_c \approx g_mv_{be}
  \quad\Longrightarrow\quad
  h_{fe}(f) = \frac{i_c}{i_b} = \frac{\beta}{1+jf/f_\beta},
\]
con
\[
  f_\beta = \frac{1}{2\pi r_\pi(C_\pi+C_\mu)} .
\]

\begin{clave}[Frecuencia de transición \enhoja]
$|h_{fe}|$ cae a 1 en
\[
  f_T = \beta\,f_\beta = \frac{g_m}{2\pi\,(C_\pi+C_\mu)}
  \;\le\; \frac{1}{2\pi\tau_F}\quad(\text{límite a corriente alta}).
\]
Por encima de $f_\beta$, $|h_{fe}|\approx f_T/f$: el producto ganancia $\times$ ancho de banda
es constante e igual a $f_T$, un dato que dan todas las hojas de datos.
\end{clave}

\begin{figure}[H]
\centering
\begin{tikzpicture}
\begin{loglogaxis}[
  width=0.75\textwidth,height=5.8cm,
  xmin=1e5,xmax=3e10,ymin=0.5,ymax=500,
  xlabel={$f$ (Hz)},ylabel={$|h_{fe}|$},
  grid=major,grid style={gray!25},
]
\addplot[azul,line width=1.3pt] table[x=f,y=hfe] {datos/hfe.dat};
\addplot[gris,dashed,domain=1e7:4e9] {2e9/x};
\draw[gris,dashed] (axis cs:1e7,0.5) -- (axis cs:1e7,141);
\draw[gris,dashed] (axis cs:2e9,0.5) -- (axis cs:2e9,1);
\node[font=\small,anchor=north west] at (axis cs:2e5,180) {$h_{fe0}=\beta=200$};
\node[font=\small,anchor=north west] at (axis cs:1.1e7,0.9) {$f_\beta$};
\node[font=\small,anchor=north west] at (axis cs:2.1e9,0.9) {$f_T$};
\node[gris,font=\scriptsize,anchor=west,rotate=-38] at (axis cs:1.5e8,16) {$-20$ dB/década};
\end{loglogaxis}
\end{tikzpicture}
\caption{Módulo de la ganancia de corriente de cortocircuito (ejemplo de la \diapo{2.33}:
$h_{FE}=200$, $f_T=\SI{2}{GHz}$, por lo que $f_\beta=\SI{10}{MHz}$).}
\label{fig:hfe}
\end{figure}

\begin{ejemplo}[$f_T$ del transistor del ejemplo]
Con $W=\SI{1}{\micro m}$ y $D_B=\SI{12.4}{cm^2/s}$: $\tau_F=W^2/(2D_B)=\SI{0.40}{ns}$. A
$I_C=\SI{1}{mA}$: $C_b=\tau_Fg_m=\SI{0.40}{ns}\times\SI{38.5}{mS}=\SI{15.4}{pF}$. Con
$C_{je0}=\SI{1}{pF}$ ($C_{je}\approx\SI{2}{pF}$) y $C_\mu=\SI{0.5}{pF}$:
\[
  f_T=\frac{\SI{38.5}{mS}}{2\pi(\SI{17.4}{pF}+\SI{0.5}{pF})}\approx\SI{340}{MHz},
  \qquad
  \frac{1}{2\pi\tau_F}\approx\SI{400}{MHz}.
\]
Para llegar a GHz hace falta una base mucho más angosta ($\tau_F\propto W^2$): con
$W=\SI{0.1}{\micro m}$, $\tau_F$ baja 100 veces.
\end{ejemplo}

\subsection{Otros modelos: parámetros Z, Y, H y G}

Un cuadripolo lineal se describe con 4 parámetros, eligiendo cuáles variables son
independientes \diapo{2.24--2.27}:
\[
  \text{Z: }\begin{bmatrix}V_1\\V_2\end{bmatrix}=\begin{bmatrix}z_{11}&z_{12}\\z_{21}&z_{22}\end{bmatrix}\begin{bmatrix}I_1\\I_2\end{bmatrix}
  \qquad
  \text{Y: }\begin{bmatrix}I_1\\I_2\end{bmatrix}=\begin{bmatrix}y_{11}&y_{12}\\y_{21}&y_{22}\end{bmatrix}\begin{bmatrix}V_1\\V_2\end{bmatrix}
  \qquad
  \text{H: }\begin{bmatrix}V_1\\I_2\end{bmatrix}=\begin{bmatrix}h_{11}&h_{12}\\h_{21}&h_{22}\end{bmatrix}\begin{bmatrix}I_1\\V_2\end{bmatrix}
\]
(G es el inverso de H.) Los parámetros \textbf{híbridos} (H) son los clásicos de las hojas de
datos de transistores, porque se miden fácilmente (cortocircuito a la salida, circuito
abierto a la entrada). En emisor común se escriben con subíndice $e$:

\begin{table}[H]
\centering\small
\caption{Parámetros $h$ en emisor común y su relación con el híbrido-$\pi$ (baja frecuencia).}
\label{tab:hparam}
\begin{tabular}{llll}
\toprule
Parámetro & Definición & Significado & Híbrido-$\pi$ \\
\midrule
$h_{ie}=h_{11}$ & $v_{be}/i_b$ con $v_{ce}=0$ & impedancia de entrada & $r_b+r_\pi\approx r_\pi$ \\
$h_{re}=h_{12}$ & $v_{be}/v_{ce}$ con $i_b=0$ & transferencia inversa de tensión & $r_\pi/r_\mu\sim10^{-4}$ \\
$h_{fe}=h_{21}$ & $i_c/i_b$ con $v_{ce}=0$ & ganancia de corriente & $g_mr_\pi=\beta$ \\
$h_{oe}=h_{22}$ & $i_c/v_{ce}$ con $i_b=0$ & admitancia de salida & $\approx1/r_o$ \\
\bottomrule
\end{tabular}
\end{table}

\begin{ojo}[Giacoletto no es el modelo híbrido \diapo{2.28, 2.30}]
Al híbrido-$\pi$ a veces se lo llama ``modelo híbrido'', pero son modelos distintos: el de
Giacoletto tiene una fuente de corriente controlada por \emph{tensión} ($g_mv_\pi$) y elementos
con significado físico ($r_\pi$, $r_o$, capacidades); el modelo H tiene una fuente controlada
por \emph{corriente} ($h_{fe}i_b$) y parámetros medidos en un punto de trabajo. La conversión
entre ambos es directa (tabla~\ref{tab:hparam}). Los parámetros $h$ dependen del punto de
trabajo y de la temperatura: los fabricantes dan curvas de corrección normalizadas
\diapo{2.31--2.32}.
\end{ojo}

\subsection{\texorpdfstring{$h_{FE}$ vs.\ $h_{fe}$}{hFE vs hfe}}

\begin{itemize}
  \item $\beta_{DC}=h_{FE}=I_C/I_B$: cociente de continuas en el punto $Q$.
  \item $\beta_{ac}=h_{fe}=\Delta i_C/\Delta i_B$: pendiente (cociente de variaciones).
\end{itemize}
Derivando $I_C=\beta_{DC}(I_C)\,I_B$ se obtiene
\[
  \beta_{ac} = \frac{\beta_{DC}}{1-\dfrac{I_C}{\beta_{DC}}\dfrac{\dd\beta_{DC}}{\dd I_C}} .
\]
Si $\beta_{DC}$ fuera constante ambos coincidirían. En la materia se toma
$\beta=\beta_{DC}=\beta_{ac}$ \diapo{2.29}.

\begin{ojo}[\diapo{2.29}]
``$\beta_{ac}$ es levemente inferior al $\beta$ de continua'' sólo vale donde $\beta_{DC}$
\emph{decrece} con $I_C$ (corrientes altas, figura~\ref{fig:gummel}). A corrientes bajas, donde
$\beta_{DC}$ crece con $I_C$, resulta $\beta_{ac}>\beta_{DC}$.
\end{ojo}

% =====================================================================
\section{Resumen y repaso}
% =====================================================================

\subsection{Las regiones de un vistazo (PNP)}

\begin{table}[H]
\centering\small
\begin{tabularx}{\textwidth}{l>{\raggedright\arraybackslash}p{2.6cm}>{\raggedright\arraybackslash}X>{\raggedright\arraybackslash}X}
\toprule
Región & Condición & Modelo & Valores típicos (Si) \\
\midrule
Activa directa & $V_{EB}>0$, $V_{CB}<0$ & $I_C=\beta I_B=I_Se^{V_{EB}/V_T}(1+V_{EC}/V_A)$ & $V_{EB}\approx\SIrange{0.6}{0.7}{V}$ \\
Saturación & $V_{EB}>0$, $V_{CB}>0$ & $I_C<\beta I_B$, fijada por el circuito & $V_{EC,sat}\approx\SIrange{0.05}{0.2}{V}$ \\
Corte & $V_{EB}<0$, $V_{CB}<0$ & $I_C\approx I_{CE0}$ o $I_{CB0}$ & $I\approx0$ \\
Activa inversa & $V_{EB}<0$, $V_{CB}>0$ & $I_E=\beta_RI_B$ & $\beta_R\sim1$--$5$ \\
\bottomrule
\end{tabularx}
\end{table}

\needspace{18\baselineskip}
\subsection{Formulario}

\begin{tcolorbox}[base,colback=white,colframe=verde,title={Formulario}]
\renewcommand{\arraystretch}{1.8}
\small
\begin{tabularx}{\linewidth}{>{\raggedright\arraybackslash}p{0.17\linewidth}X}
Base angosta\newline\enhojatexto & $\Delta p_B(x)=\Delta p_B(0)\left(1-\frac{x}{W}\right)+\Delta p_B(W)\frac{x}{W}$,\quad $\Delta p_B(0)=p_{B0}\eVEB$ \\
Corrientes\newline\enhojatexto & $I_{En}=qA\frac{D_E}{L_E}n_{E0}\eVEB$,\quad $I_{Cn}=-qA\frac{D_C}{L_C}n_{C0}\eVCB$ \\
 & $I_{Ep}$, $I_{Cp}$ con $\tanh$ y $\sinh$ (sección~\ref{sec:formal});\quad $I_B=I_{En}+I_{Cn}+I_{RB}$ \\
Parámetros\newline\enhojatexto & $\gamma=\frac{I_{Ep}}{I_E}$,\quad $\aT=\frac{I_{Cp}}{I_{Ep}}\approx\frac{1}{1+\frac12(W/L_B)^2}$,\quad $\gamma\approx\frac{1}{1+\frac{D_E}{D_B}\frac{N_B}{N_E}\frac{W}{L_E}}$ \\
 & $\alpha=\gamma\aT$,\quad $\beta=\frac{\alpha}{1-\alpha}\approx\left[\frac{D_E}{D_B}\frac{N_B}{N_E}\frac{W}{L_E}+\frac12\left(\frac{W}{L_B}\right)^2\right]^{-1}$ \\
Ebers--Moll\newline\enhojatexto & $I_E=I_{F0}\eVEB-\aR I_{R0}\eVCB$ \\
 & $I_C=\aF I_{F0}\eVEB-I_{R0}\eVCB$,\quad $\aF I_{F0}=\aR I_{R0}$ \\
Fugas\newline\enhojatexto & $I_C=\alpha I_E+I_{CB0}=\beta I_B+I_{CE0}$,\quad $I_{CE0}=(\beta+1)I_{CB0}$ \\
Ancho de base & $W=W_B-x_{nEB}-x_{nCB}$ \\
Early\newline\enhojatexto & $r_o=V_A/I_C$,\quad $V_A\approx W/|\dd W/\dd V_{CB}|$ \\
Avalancha & $M=\left[1-(V_{CB}/BV_{CB0})^n\right]^{-1}$,\quad $BV_{CE0}\approx BV_{CB0}\,\beta^{-1/n}$ \\
Saturación & $V_{EC,sat}=V_T\ln\frac{1/\aR+\beta_{forz}/\beta_R}{1-\beta_{forz}/\beta_F}$,\quad $\beta_{forz}=I_C/I_B<\beta$ \\
Pequeña señal\newline\enhojatexto & $g_m=\frac{I_C}{V_T}$,\quad $r_\pi=\frac{\beta}{g_m}$,\quad $r_o=\frac{V_A}{I_C}$,\quad $r_\mu\approx\beta r_o$,\quad $g_mr_o=\frac{V_A}{V_T}$ \\
Capacidades\newline\enhojatexto & $C_\pi=C_b+C_{je}$,\quad $C_b=\tau_Fg_m$ ($\tau_F=\frac{W^2}{2D_B}$),\quad $C_{je}\approx2C_{je0}$,\quad $C_\mu=\frac{C_{\mu0}}{(1-V_{CB}/V_{bi})^{n_c}}$ \\
Frecuencia\newline\enhojatexto & $f_\beta=\frac{1}{2\pi r_\pi(C_\pi+C_\mu)}$,\quad $f_T=\beta f_\beta=\frac{g_m}{2\pi(C_\pi+C_\mu)}$ \\
\end{tabularx}
\end{tcolorbox}

\subsection{Preguntas de repaso}

\begin{enumerate}
  \item \textbf{¿Por qué la base tiene que ser angosta y poco dopada respecto del emisor?}\\
  Angosta ($W\ll L_B$) para que los portadores inyectados la crucen sin recombinarse
  ($\aT\to1$); menos dopada que el emisor para que la juntura E--B inyecte casi sólo hacia la
  base ($\gamma\to1$).
  \item \textbf{¿Por qué $I_C$ no depende (casi) de $V_{CB}$ en activa?}\\
  Porque la juntura C--B en inversa sólo ``barre'' lo que le llega; lo que llega lo fija la
  inyección de E--B ($e^{V_{EB}/V_T}$). La dependencia residual es el efecto Early.
  \item \textbf{¿Qué parámetro limita $\beta$ en un transistor moderno?}\\
  La eficiencia de emisión: $1-\gamma\propto W/L_E$ suele ser mucho mayor que
  $1-\aT\propto(W/L_B)^2$.
  \item \textbf{¿Por qué un BJT usado al revés (colector como emisor) es tan malo?}\\
  Porque el colector está poco dopado: la eficiencia de emisión inversa es baja
  ($\aR\sim0.5$--$0.8$, $\beta_R\sim1$--$5$).
  \item \textbf{¿Por qué $BV_{CE0}<BV_{CB0}$ e $I_{CE0}>I_{CB0}$?}\\
  Con la base abierta, la corriente generada en C--B actúa como corriente de base y el
  transistor la amplifica por $\beta+1$.
  \item \textbf{¿Cómo puede ser $V_{EC}\approx\SI{0.1}{V}$ con $V_{EB}=\SI{0.7}{V}$?}\\
  En saturación $V_{CB}$ también es positiva y casi igual: $V_{EC}=V_{EB}-V_{CB}$.
  \item \textbf{¿De qué depende $g_m$?}\\
  Sólo de la corriente de polarización: $g_m=I_C/V_T$. No depende de $\beta$, del área ni del
  material.
  \item \textbf{¿Qué fija la frecuencia de transición?}\\
  La carga que hay que mover para cambiar la corriente: el tiempo de tránsito por la base
  ($\tau_F=W^2/2D_B$) y las capacidades de juntura. A corriente alta $f_T\to1/(2\pi\tau_F)$.
  \item \textbf{¿Por qué un transistor en saturación tarda en apagarse?}\\
  Porque tiene mucha carga de minoritarios almacenada en la base y el colector, que hay que
  retirar antes de que la corriente se corte.
\end{enumerate}

\subsection{Correcciones a las diapositivas (resumen)}

\begin{table}[H]
\centering\small
\begin{tabularx}{\textwidth}{cX}
\toprule
Diapo & Corrección \\
\midrule
1.14 & $I_{Cn}=+qAD_C\,\dd\Delta n_C/\dd x'$ (el resultado final de la diapo es correcto). \\
1.15 & Las condiciones de contorno son de la \emph{base}, no del colector. \\
1.16 & $I_{Ep}=-qAD_B\,\dd\Delta p_B/\dd x|_0$ (falta el signo de la corriente de difusión); ídem $I_{Cp}$. \\
1.28 & Las curvas $i_C$ vs.\ $v_{BE}$ son de transferencia; la de entrada es $i_B$ vs.\ $v_{BE}$. \\
1.31, 2.34 & ``Controlado por corriente'' es un modelo de circuito; físicamente $I_C$ lo controla $V_{EB}$. \\
2.3 & En un PNP en MAD la condición es $V_{CB}<-5V_T$, no $V_{CB}>5V_T$. \\
2.8 & $BV_{CE0}=0.52\,BV_{CB0}$ es un caso particular; en general $BV_{CE0}\approx BV_{CB0}\beta^{-1/n}$. \\
2.11 & $I_{CB0}=I_{R0}(1-\aF\aR)$ (más generación en la ZCE), no sólo $I_{Cn}$. \\
2.21 & $C_{De}$ con $L_B$ es de diodo largo; en la base angosta $C_b=\tau_Fg_m$. El $\tau_B$ de la hoja es tiempo de tránsito, no vida media. \\
2.29 & $\beta_{ac}$ puede ser mayor o menor que $\beta_{DC}$ según la pendiente de $\beta_{DC}(I_C)$. \\
\bottomrule
\end{tabularx}
\end{table}

\subsection*{Bibliografía}
\begin{itemize}[nosep]
  \item R.~F.~Pierret, \emph{Semiconductor Device Fundamentals}, Addison-Wesley, 1996: caps.~10
  (fundamentos), 11 (características estáticas) y 12 (respuesta dinámica).
  \item A.~S.~Sedra, K.~C.~Smith, \emph{Microelectronic Circuits}, 7.\textsuperscript{a} ed., Oxford, 2015: cap.~6.
  \item Diapositivas de clase ``(25.18) Física Electrónica -- Transistor bipolar de juntura'',
  1.\textsuperscript{a} y 2.\textsuperscript{a} parte, ITBA.
  \item Hoja de fórmulas y tabla de constantes de la materia.
\end{itemize}


\end{document}
